% ===========================================================================
% APPENDIX A: PROOFS AND AUXILIARY MATERIAL
% (self-contained restatement of assumptions and results; statements are
%  identical to Section 4 of the manuscript; numbering prefixed A)
% ===========================================================================
\appendix
\setcounter{theorem}{0}\setcounter{lemma}{0}\setcounter{corollary}{0}
\setcounter{assumption}{0}\setcounter{remark}{0}\setcounter{equation}{0}
\renewcommand{\thetheorem}{A\arabic{theorem}}
\renewcommand{\thelemma}{A\arabic{lemma}}
\renewcommand{\thecorollary}{A\arabic{corollary}}
\renewcommand{\theassumption}{A\arabic{assumption}}
\renewcommand{\theremark}{A\arabic{remark}}
\renewcommand{\theequation}{A\arabic{equation}}
\renewcommand{\thefigure}{A\arabic{figure}}

\section{Proofs and auxiliary material}\label{app:proofs}

This appendix contains the complete proofs of all results stated in
Section~4 of the main text, together with the auxiliary lemmas used in
them. Assumptions, lemmas, and theorems are restated for convenience;
statements are identical to the main text, and numbering is prefixed
A.

\subsubsection{Notation and setup}\label{app:sec:setup}

We adopt the notation of the manuscript. For $i=1,\dots,n$, let
$Y_i\in\N_0$ be a count outcome and $X_i\in\R^p$ an unobserved covariate,
with
\[
  \E[Y_i\mid X_i]=\exp(X_i^\top\bet_0),
\]
where $\bet_0\in\Theta\subset\R^p$ is the parameter of interest. The observed
covariate is $W_i=X_i+\beps_i$, where $\beps_i\sim N(\bm 0,\bSig_\epsilon)$
is independent of $(X_i,Y_i)$, with $\bSig_\epsilon=\operatorname{diag}
(\sigma^2_{\epsilon,1},\dots,\sigma^2_{\epsilon,p})$ positive semidefinite
and known. Define the augmented covariate, for $\lambda\ge 0$ and a
simulated error draw $U_i\sim N(\bm 0,\bSig_\epsilon)$ independent of
everything else,
\[
  W_i(\lambda) \;=\; W_i+\sqrt{\lambda}\,U_i \;=\; X_i+\beps_i+\sqrt{\lambda}\,U_i .
\]
For a grid $\Lambda=\{\lambda_0,\dots,\lambda_{K-1}\}$ with
$0=\lambda_0<\lambda_1<\cdots<\lambda_{K-1}=\lambda_{\max}$, the SIMEX
estimator at level $\lambda_j$ is
\[
  \bar\beta(\lambda_j)=\frac1B\sum_{b=1}^{B}\hat\beta_b(\lambda_j),\qquad
  \hat\beta_b(\lambda_j)=\argmax_{\bet\in\Theta}\;
  \frac1n\sum_{i=1}^{n}\Bigl[Y_i\,W_i^{(b)}(\lambda_j)^\top\bet
  -\exp\bigl(W_i^{(b)}(\lambda_j)^\top\bet\bigr)\Bigr],
\]
where $W_i^{(b)}(\lambda_j)=W_i+\sqrt{\lambda_j}\,U_{ib}$ and
$\{U_{ib}\}$ are i.i.d.\ $N(\bm 0,\bSig_\epsilon)$ across both $i$ and $b$,
independent of the data. The \asimex{} estimator fits a degree-$d$
polynomial $\hat p(\lambda)=\bphi(\lambda)^\top\hat\theta$ by least squares
to the $K$ points $(\lambda_j,\bar\beta(\lambda_j))$, where
$\bphi(\lambda)=(1,\lambda,\dots,\lambda^d)^\top$, and extrapolates:
\[
  \hat\beta_{\asimex} \;=\; \hat p(-1)\;=\;\bphi(-1)^\top\hat\theta,
  \qquad
  \hat\theta=\argmin_{\theta\in\R^{(d+1)p}}
  \sum_{j=0}^{K-1}\bigl\|\bar\beta(\lambda_j)-\bphi(\lambda_j)^\top\theta\bigr\|^2 .
\]
Throughout, $\|\cdot\|$ denotes the Euclidean norm, $\Rightarrow$ weak
convergence in $\ell^\infty$, $\rightarrow_p$ convergence in probability and
$\rightarrow_d$ convergence in distribution. Generic positive constants are
denoted $C,c$ and may change from line to line.

% ===========================================================================
\subsubsection{Corrected assumptions}\label{app:sec:assumptions}

The assumption set below replaces Assumptions~1--5 of the draft. Two changes
are substantive: the cross-validation assumption (A3) is restated as an
explicit \emph{selection-stability} condition (the draft's (A3b), that the
selected pair converges ``with positive probability,'' is not sufficient for
the proofs); and a new assumption~(A6) quantifies the polynomial
extrapolation bias, replacing the incorrect inequality~(43) of the draft.

\begin{assumption}[Model and regularity]\label{app:ass:model}
\leavevmode
\begin{enumerate}\itemsep2pt
\item[(a)] $Y_1,X_1,\beps_1,\dots,Y_n,X_n,\beps_n$ are i.i.d.;
      $\E[Y_i\mid X_i]=\exp(X_i^\top\bet_0)$ with $\bet_0$ in the interior
      of the compact convex set $\Theta\subset\R^p$.
\item[(b)] $\E\|X_i\|^4<\infty$, $\E|Y_i|^2<\infty$, and
      $\E[X_iX_i^\top]$ is positive definite.
\item[(c)] $\beps_i\sim N(\bm 0,\bSig_\epsilon)$ with $\bSig_\epsilon$
      known, positive semidefinite, independent of $(X_i,Y_i)$.
\end{enumerate}
\end{assumption}

\begin{assumption}[Pseudo-design non-collinearity]\label{app:ass:pseudo}
For every $\delta$ in a neighbourhood of $[-1,\lambda_{\max}+1]$ and every
$\bet$ in a neighbourhood of $\bet_0$,
\[
  \bm H(\bet,\delta)
  := \E\Bigl[\exp\!\bigl(\bet^\top(X_i+\beps_\delta)\bigr)
      (X_i+\beps_\delta)(X_i+\beps_\delta)^\top\Bigr]
\]
is positive definite, where $\beps_\delta\sim N(\bm 0,\delta\bSig_\epsilon)$
and the expectation is over $(X_i,Y_i,\beps_\delta)$. In particular
$\inf_{\bet,\delta}\lambda_{\min}\bigl(\bm H(\bet,\delta)\bigr)\ge c_H>0$.
\end{assumption}

\begin{assumption}[Simulation and grid]\label{app:ass:simgrid}
The number of simulation replicates satisfies $B=B(n)\to\infty$. The grid
$\Lambda_n$ has $K=K(n)$ points, equally spaced on $[0,\lambda_{\max}]$,
with $K\to\infty$, $\lambda_{\max}\le\bar\lambda<\infty$ fixed, and
$K^2\log n=o(n)$.
\end{assumption}

\begin{assumption}[Extrapolation degree]\label{app:ass:degree}
The degree grid is $\{1,\dots,d_{\max}(n)\}$ with $d_{\max}(n)\to\infty$,
$d_{\max}(n)\le K(n)/4$, and the selected degree
$\hat d_n\in\{1,\dots,d_{\max}(n)\}$ satisfies
$\hat d_n\to\infty$ in probability and
$d_{\max}(n)^3/K(n)\to 0$.
\end{assumption}

\begin{assumption}[Selection stability, for the central limit
theorem]\label{app:ass:cv}
The cross-validation procedure selects a pair
$(\hat\lambda_{\max,n},\hat d_n)$ from a fixed finite grid
$\mathcal G$ such that
$\Prob\bigl((\hat\lambda_{\max,n},\hat d_n)=(\lambda^*,d^*)\bigr)\to 1$
for some $(\lambda^*,d^*)\in\mathcal G$.
\end{assumption}

\begin{remark}[Why (A4)--(A6) take this form]
\leavevmode
\begin{enumerate}\itemsep2pt
\item[(i)] The draft's Assumption~4(a) imposed $B/n\to 0$ \emph{together
with} $B/\sqrt n\to\infty$. As shown in the proof of Theorem~\ref{app:thm:an}
below, the simulation component of the \asimex{} estimator is
$O_p(B^{-1/2})$ at the $\sqrt n$ scale, so \emph{any} $B\to\infty$
suffices for the CLT; conversely $B/\sqrt n\to\infty$ is neither necessary
nor (with $B/n\to 0$) the natural pairing. We require only $B\to\infty$.
\item[(ii)] The draft's inequality~(43) claimed
$\max_{\lambda\in[0,\lambda_{\max}]}|p(\lambda)-\beta(\lambda)|
\le C\lambda_{\max}^{-p^*_{\deg}}\to 0$ ``as $\lambda_{\max}$ grows'' for
\emph{fixed} polynomial degree. This is false: for fixed degree the
polynomial approximation error on a compact interval does not vanish as the
interval endpoint grows. The corrected route is to let the CV-selected
degree grow with $n$ (Assumption~\ref{app:ass:degree}) while the real-analyticity
of the population curve (Lemma~\ref{app:lem:curve} below) makes the
degree-$d$ best-approximation error decay geometrically, so that the
extrapolation bias is $o(n^{-1/2})$; see Lemma~\ref{app:lem:approx}.
\item[(iii)] Assumption~\ref{app:ass:cv} is used only for Theorem~\ref{app:thm:an}.
It is mild: the hyperparameter grid is finite and the CV risk converges
uniformly over it, so the selected pair is eventually constant with
probability tending to one; the consistency theorem
(Theorem~\ref{app:thm:cons}) instead uses the growing-degree regime of
Assumption~\ref{app:ass:degree}. A fully data-driven treatment reconciling
(A4) and (A5) is possible but requires a model-selection argument
(e.g.\ penalizing the CV risk by $d\log K/n$); we relegate it to future
work and note it in Remark~\ref{app:rem:cvreconcile}.
\end{enumerate}
\end{remark}

% ===========================================================================
\subsubsection{The population SIMEX curve}\label{app:sec:population}

The proofs below rest on a structural lemma that replaces the incorrect
Step~2 of the draft's consistency proof. Recall
$\delta=1+\lambda$; the augmented error $\beps_i+\sqrt\lambda U_i$ is
$N(\bm 0,\delta\bSig_\epsilon)$.

\begin{lemma}[Population representation and analyticity]\label{app:lem:curve}
Grant Assumptions~\ref{app:ass:model} and~\ref{app:ass:pseudo}. Define, for
$\delta\ge 0$,
\begin{equation}\label{app:eq:popobj}
  Q(\bet;\delta)
  \;=\; \E\Bigl[Y_i\,\bet^\top(X_i+\beps_\delta)
      -\exp\bigl(\bet^\top(X_i+\beps_\delta)\bigr)\Bigr],
  \qquad
  \beta(\delta)=\argmax_{\bet\in\Theta}Q(\bet;\delta),
\end{equation}
and let $\beta(\lambda):=\beta(1+\lambda)$. Then:
\begin{enumerate}\itemsep2pt
\item[(i)] $Q$ depends on $\delta$ only through the scalar function
      \[
        Q(\bet;\delta)\;=\;\E\bigl[Y_iX_i^\top\bet\bigr]
        \;-\;\E\Bigl[\exp\bigl(X_i^\top\bet\bigr)\Bigr]\,
        \exp\!\Bigl(\tfrac{\delta}{2}\,\bet^\top\bSig_\epsilon\bet\Bigr),
      \]
      hence $Q(\cdot\,;\delta)$ is well defined and real-analytic in
      $(\bet,\delta)$ on $\operatorname{int}\Theta\times(-1-\eta,\bar\lambda+1)$
      for some $\eta>0$ (the Gaussian exponential tilt defines $Q$ for
      negative $\delta$ as well).
\item[(ii)] The maximizer $\beta(\delta)$ exists, is unique, lies in the
      interior of $\Theta$, and is real-analytic in $\delta$ on a
      neighbourhood of $[-1,\bar\lambda+1]$.
\item[(iii)] $\beta(-1)=\beta(0)=\bet_0$; moreover
      $\beta(\delta)=\bet_0+\delta\,\tilde{\bm b}+O(\delta^2)$ as
      $\delta\downarrow 0$, where, writing
      $\bm g(\bet)=\E\bigl[e^{X_i^\top\bet}X_i\bigr]$ and
      $M(\bet)=\E\bigl[e^{X_i^\top\bet}\bigr]$,
      \begin{equation}\label{app:eq:bias1}
        \tilde{\bm b}
        =-\bm H(\bet_0,0)^{-1}\Bigl[
          \tfrac12\bigl(\bet_0^\top\bSig_\epsilon\bet_0\bigr)\,
          \bm g(\bet_0)\;+\;\bSig_\epsilon\bet_0\,M(\bet_0)
        \Bigr]
      \end{equation}
      is the leading-order bias of the naive estimator (since
      $\beta(1)-\bet_0=\tilde{\bm b}+O(\|\bSig_\epsilon\|^2)$). For
      $p=1$ and $X\sim N(0,1)$ this gives
      $\tilde b=-\sigma^2\beta_0\bigl(1+\tfrac12\beta_0^2\bigr)
      /(1+\beta_0^2)<0$ for $\beta_0>0$, the familiar attenuation.
\end{enumerate}
\end{lemma}

\begin{proof}
\emph{(i)} Since $\beps_\delta=\beps_i+\sqrt\lambda U_i$ with
$\beps_i,U_i$ independent $N(\bm 0,\bSig_\epsilon)$, we have
$\beps_\delta\sim N(\bm 0,(1+\lambda)\bSig_\epsilon)$ and
$\beps_\delta\perp(X_i,Y_i)$. For Gaussian $\beps_\delta$, the exponential
tilting identity gives, for any $\bet$,
\[
  \E\Bigl[\exp\bigl(\bet^\top(X_i+\beps_\delta)\bigr)\Bigm|X_i\Bigr]
  =\exp\bigl(\bet^\top X_i\bigr)\,
   \exp\!\Bigl(\tfrac{\delta}{2}\bet^\top\bSig_\epsilon\bet\Bigr),
\]
and $\E[\bet^\top(X_i+\beps_\delta)\mid X_i,Y_i]=\bet^\top X_i$ because
$\E[\beps_\delta]=0$ and $\beps_\delta\perp(X_i,Y_i)$. Taking expectations
yields the displayed factorization. Analyticity of
$\exp(X^\top\bet)\exp(\delta\bet^\top\bSig_\epsilon\bet/2)$ in
$(\bet,\delta)$ is immediate, and the expectation is finite and
locally uniformly bounded because $X$ has finite fourth moment and
$\delta$ ranges over a bounded interval (the integrand is dominated by
$\exp(C\|X\|)$-type functions integrated against the Gaussian tilt, which
has all moments finite). This proves (i).

\emph{(ii)} Fix $\delta$ in the stated range. The Hessian of $-Q$ in
$\bet$ is
\[
  -\nabla_{\bet}^2Q(\bet;\delta)
  =\E\Bigl[\exp\bigl(\bet^\top(X_i+\beps_\delta)\bigr)
   (X_i+\beps_\delta)(X_i+\beps_\delta)^\top\Bigr]
  =\bm H(\bet,\delta),
\]
which is positive definite by Assumption~\ref{app:ass:pseudo}; hence $Q$ is
strictly concave and its maximizer over the convex compact $\Theta$ is
unique. At $\bet_0$ (which is interior by Assumption~\ref{app:ass:model}(a))
the score
$\nabla_{\bet} Q(\bet_0;0)=\E[Y_iX_i]-\E[e^{X_i^\top\bet_0}X_i]=\bm 0$
by the first-order condition of the true Poisson model, so for $\delta$
in a neighbourhood of $0$ the maximizer lies in the interior of $\Theta$
for all $\delta$ in a neighbourhood of $[-1,\bar\lambda+1]$
(by continuity of the score along the compact segment and
$\lambda_{\min}\bm H\ge c_H$; a margin argument gives a uniform interiority
bound). The implicit function theorem applied to
$\nabla_{\bet} Q(\bet;\delta)=\bm 0$ at $(\bet,\delta)=(\beta(\delta),\delta)$,
with Jacobian $-\bm H(\bet,\delta)\succ 0$ invertible and
$Q$ real-analytic by (i), yields that $\delta\mapsto\beta(\delta)$ is
real-analytic on the stated neighbourhood. This proves (ii).

\emph{(iii)} At $\delta=0$ the objective
$Q(\bet;0)=\E[Y_iX_i^\top\bet]-\E[e^{X_i^\top\bet}]$ is the population
Poisson log-likelihood with covariate $X_i$, which by the Kullback--Leibler
argument is uniquely maximized at $\bet_0$: indeed
$Q(\bet;0)-Q(\bet_0;0)=\E[e^{X_i^\top\bet_0}(X_i^\top(\bet-\bet_0))
- e^{X_i^\top(\bet-\bet_0)}+1]\le 0$ with equality iff
$X_i^\top(\bet-\bet_0)=0$ almost surely, i.e.\ $\bet=\bet_0$ by
Assumption~\ref{app:ass:model}(b). Hence $\beta(0)=\bet_0$ and, since
$\beta(-1)=\beta(\delta=0)$, also $\beta(-1)=\bet_0$.

For the expansion, differentiate the implicit relation
$\nabla_{\bet} Q(\beta(\delta);\delta)=\bm 0$ at $\delta=0$. By the implicit
function theorem, $\beta'(\delta)=\bm H(\beta(\delta),\delta)^{-1}\,
\partial_\delta\nabla_{\bet} Q\,(\beta(\delta);\delta)$. From the
factorization in (i) together with the two Gaussian-tilt identities
\[
  \E\bigl[e^{\bet^\top(X_i+\beps_\delta)}\bigr]
  =e^{\frac{\delta}{2}\bet^\top\bSig_\epsilon\bet}M(\bet),
  \qquad
  \E\bigl[e^{\bet^\top(X_i+\beps_\delta)}(X_i+\beps_\delta)\bigr]
  =e^{\frac{\delta}{2}\bet^\top\bSig_\epsilon\bet}
   \bigl(\bm g(\bet)+\delta\,\bSig_\epsilon\bet\,M(\bet)\bigr),
\]
(the second follows from
$\E[e^{\bet^\top\beps_\delta}\beps_\delta]
=\delta\,\bSig_\epsilon\bet\,e^{\frac{\delta}{2}\bet^\top\bSig_\epsilon\bet}$
for $\beps_\delta\sim N(\bm 0,\delta\bSig_\epsilon)$), the population score
is
\[
  \nabla_{\bet} Q(\bet;\delta)=\E[Y_iX_i]
  -e^{\frac{\delta}{2}\bet^\top\bSig_\epsilon\bet}
  \bigl(\bm g(\bet)+\delta\,\bSig_\epsilon\bet\,M(\bet)\bigr).
\]
Differentiating in $\delta$ at $\delta=0$,
\[
  \partial_\delta\nabla_{\bet} Q(\bet;\delta)\big|_{\delta=0}
  =-\tfrac12\bigl(\bet^\top\bSig_\epsilon\bet\bigr)\bm g(\bet)
   -\bSig_\epsilon\bet\,M(\bet),
\]
so at $\bet=\bet_0$,
$\beta'(0)=\bm H(\bet_0,0)^{-1}\partial_\delta\nabla_{\bet} Q(\bet_0;0)
=\tilde{\bm b}$ as displayed in \eqref{app:eq:bias1}. The scalar normal check:
$\mu_0(\bet)=\E[e^{\bet X}]=e^{\bet^2/2}$,
$\mu_1(\bet)=\bet e^{\bet^2/2}$,
$\mu_2(\bet)=(1+\bet^2)e^{\bet^2/2}$, so
$\tilde b=-\bigl[\tfrac12\beta_0^3+\beta_0\bigr]e^{\beta_0^2/2}
/\bigl((1+\beta_0^2)e^{\beta_0^2/2}\bigr)
=-\sigma^2$-free multiple $-\beta_0(1+\beta_0^2/2)/(1+\beta_0^2)$
times $\sigma^2$ carried by $\bSig_\epsilon$; this matches the direct
computation from the scalar pseudo-score equation
$e^{\frac{\delta}{2}\beta^2\sigma^2}\beta(1+\delta\sigma^2)e^{\beta^2/2}
=\beta_0e^{\beta_0^2/2}$, confirming attenuation. Real-analyticity from
(ii) then gives $\beta(\delta)=\bet_0+\delta\tilde{\bm b}+O(\delta^2)$.
This proves (iii).
\end{proof}

\begin{remark}
Lemma~\ref{app:lem:curve} is the correct replacement for draft Step~2, which
argued that ``$W(-1)=W-\epsilon=X$''. That identity is invalid: the
simulated noise $U$ is \emph{independent} of the observed error $\beps$,
so $W-\sqrt{-1}\,U=W-U\neq X$; the observed $\beps$ can never be
subtracted off. The correct statement is distributional: at level
$\lambda$, the augmented covariate is $X+\beps+\sqrt\lambda U$ with
$(\beps,U)$ i.i.d.\ $N(0,\bSig_\epsilon)$, i.e.\ $X+\beps_\delta$ with
$\delta=1+\lambda$, and the pseudo-true parameter $\beta(\delta)$ depends
on $\delta$ only through the total error variance. Setting $\delta=0$
($\lambda=-1$) removes \emph{all} error variance from the pseudo-population,
which is why $\beta(-1)=\bet_0$ holds exactly even though no single
realization ever equals $X$.
\end{remark}

% ===========================================================================
\subsubsection{Uniform convergence of the empirical SIMEX curve}\label{app:sec:curve}

Let
\[
  Q_n(\bet;\lambda,b)=\frac1n\sum_{i=1}^n
  \Bigl[Y_i\,W_i^{(b)}(\lambda)^\top\bet
  -\exp\bigl(W_i^{(b)}(\lambda)^\top\bet\bigr)\Bigr],
\]
and note that
$\E\bigl[Q_n(\bet;\lambda,b)\bigr]=Q(\bet;1+\lambda)$ for every
$(\lambda,b)$, by Lemma~\ref{app:lem:curve}(i). Write
$\beta(\lambda):=\beta(1+\lambda)$ and
$\bm H(\lambda):=\bm H(\beta(\lambda),1+\lambda)$.

\begin{lemma}[Uniform M-estimation over grid and replicates]
\label{app:lem:unif}
Grant Assumptions~\ref{app:ass:model}--\ref{app:ass:simgrid}. Then
\[
  \max_{0\le j<K}\ \max_{1\le b\le B}
  \sup_{\bet\in\Theta}\bigl|Q_n(\bet;\lambda_j,b)-Q(\bet;1+\lambda_j)\bigr|
  \;\xrightarrow{p}\;0,
\]
and consequently
\[
  \max_{0\le j<K}\ \|\bar\beta(\lambda_j)-\beta(\lambda_j)\|
  =O_p(n^{-1/2}),
\]
with each $\hat\beta_b(\lambda_j)$ well defined and unique with probability
tending to one.
\end{lemma}

\begin{proof}
\emph{Step 1: envelope and Lipschitz bounds.} On $\Theta$ (compact) and
$\lambda\in[0,\bar\lambda]$, for any $\bet_1,\bet_2\in\Theta$,
\[
  \bigl|q(y,w;\bet_1)-q(y,w;\bet_2)\bigr|
  \le \bigl(1+|y|+\|w\|\bigr)e^{\bar\lambda\,\bm 1^\top\bSig_\epsilon\bm 1/2
        + C_{\Theta}\|w\|}\,\|\bet_1-\bet_2\|,
\]
where $q(y,w;\bet)=y\bet^\top w-e^{\bet^\top w}$ and $C_\Theta=
\sup_{\bet\in\Theta}\|\bet\|$: indeed, pointwise,
\[
  |q(y,w;\bet_1)-q(y,w;\bet_2)|\le \bigl(|y|\|w\|
  + e^{\bet_1^\top w}\|w\| + e^{\bet_2^\top w}\|w\|\bigr)\,
  \|\bet_1-\bet_2\|,
\]
and
$\bet^\top w\le\|\bet\|\|w\|$ while
$W_i=X_i+\beps_i+\sqrt\lambda U_i$ has a Gaussian tilt: under
$\Prob(\cdot\mid X_i)$, $\|W_i\|\le\|X_i\|+\|\beps_i+\sqrt\lambda U_i\|$
with the Gaussian part sub-Gaussian, giving
$\E[e^{C\|W_i\|}\mid X_i]\le e^{C\|X_i\|+C'}$. Since $\E\|X_i\|^4<\infty$,
the envelope $F(w,y)=(1+|y|+\|w\|)e^{C_{\Theta}\|w\|}$ satisfies
$\E F(W_i,Y_i)^2<\infty$ and exponential moment bounds of the
form
$\E\exp\bigl(tF(W_i,Y_i)\bigr)<\infty$ for some $t>0$ (Gaussian
$\beps_i$ supplies the sub-exponential tail of $e^{C_\Theta\|W_i\|}$).
The same argument bounds $\E F^2(W_i^{(b)}(\lambda),Y_i)$ uniformly in
$(\lambda,b)$, with constants independent of $b$ because
$W_i^{(b)}(\lambda)\stackrel{d}{=}W_i(\lambda)$ for every $b$.

\emph{Step 2: pointwise uniform LLN conditional on the simulated errors.}
Fix $(\lambda,b)$ and condition on the simulated draws
$\{U_{ib}\}_{i=1}^n$. Conditional on these, the summands
$q(Y_i,W_i^{(b)}(\lambda);\bet)$ are still i.i.d.\ across $i$? Not
conditionally (the $U_{ib}$ are fixed conditioning variables), but
unconditionally
$\{(Y_i,X_i,\beps_i,U_{ib})\}_{i=1}^n$ are i.i.d.\ as $(p+2p+1)$-tuples
indexed by $b$, so $Q_n(\bet;\lambda,b)$ is an average of i.i.d.\
mean-$Q(\bet;1+\lambda)$ summands for each $(\lambda,b)$. Standard
M-estimation uniform convergence with a sub-exponential envelope
(e.g.\ Lemma~2.4 and Theorem~2.7 of van der Vaart and Wellner, or
Bernstein's inequality for the class
$\mathcal Q_{\lambda,b}=\{q(\cdot,\cdot;\bet):\bet\in\Theta\}$, which has
finite pseudo-dimension $p+1$ and hence bounded entropy
$\log N(\epsilon\|F\|,\mathcal Q_{\lambda,b},\|\cdot\|_{L_2})
\lesssim p\log(1/\epsilon)$) gives, for each fixed $(\lambda,b)$,
\[
  \Prob\Bigl(\sup_{\bet\in\Theta}
  \bigl|Q_n(\bet;\lambda,b)-Q(\bet;1+\lambda)\bigr|>\varepsilon\Bigr)
  \le C\exp\bigl(-c\,n\varepsilon^2\bigr)
\]
for all $\varepsilon\ge C'\sqrt{\log n/n}$, with constants uniform in
$(\lambda,b)$ by the uniformity of the envelope bound in Step~1.

\emph{Step 3: union over grid and replicates.} The pairs
$(\lambda_j,b)$, $j=0,\dots,K-1$, $b=1,\dots,B$, are at most $KB$ in
number, but the corresponding summands are \emph{not} independent across
$b$ (they share the data). We therefore argue conditionally on
$\bm U=\{U_{ib}\}$: conditional on $\bm U$, for each $j$ the $B$
statistics $\sup_{\bet}|Q_n(\bet;\lambda_j,b)-Q(\bet;1+\lambda_j)|$,
$b=1,\dots,B$, are functions of the \emph{same} data with different,
independent simulated perturbations; treat instead the class
\[
  \mathcal Q_n=\bigl\{q(y,x+\beps+\sqrt{\lambda_j}u;\bet):
  \bet\in\Theta,\ j=0,\dots,K-1\bigr\},
\]
which has finite pseudo-dimension $(p+2)$ and entropy uniform in $K$
(the $\lambda_j$ enter only through the scalar $\sqrt{\lambda_j}\le
\sqrt{\bar\lambda}$, and $\{(x,\beps,u)\mapsto
\sqrt{\lambda}u\}$ is a one-dimensional parametric family Lipschitz in
$\sqrt\lambda$). Applying the maximal inequality of Theorem~2.14.1 of van
der Vaart and Wellner to the i.i.d.\ tuples
$(Y_i,X_i,\beps_i,U_{i1},\dots,U_{iB})$ over the class
$\mathcal Q_n\times\{1,\dots,B\}$ (i.e.\ augmenting the class by the
replicate index, which only multiplies the entropy by $\log B$) yields
\[
  \E\max_{j,b}\sup_{\bet\in\Theta}\bigl|Q_n(\bet;\lambda_j,b)
  -Q(\bet;1+\lambda_j)\bigr|
  \;\le\; C\sqrt{\frac{p\log(KB)+\log n}{n}},
\]
where the extra $\log n$ covers the exponential-moment argument of the
sub-exponential envelope. Since $KB\le K n$ (as $B\le n$ can be assumed
without loss; if $B>n$ replicate blocks of size $n$ and average), we get
\[
  \max_{j,b}\sup_{\bet\in\Theta}\bigl|Q_n(\bet;\lambda_j,b)
  -Q(\bet;1+\lambda_j)\bigr|
  =O_p\!\Bigl(\sqrt{\frac{\log(KB n)}{n}}\Bigr)\xrightarrow{p}0
\]
under Assumption~\ref{app:ass:simgrid} (indeed under much weaker conditions
on $K$ and $B$). This is the first claim.

\emph{Step 4: argmax convergence.} By Lemma~\ref{app:lem:curve}(ii), each
$Q(\cdot\,;1+\lambda_j)$ has a unique well-separated maximizer
$\beta(\lambda_j)$ in the interior of $\Theta$ with
$\lambda_{\min}\bm H(\lambda_j)\ge c_H$; a standard identifiability
argument (there is $\alpha>0$ with
$\sup_{\|\bet-\beta(\lambda_j)\|\ge\alpha}
\bigl(Q(\beta(\lambda_j);1+\lambda_j)-Q(\bet;1+\lambda_j)\bigr)\ge\delta_\alpha$
uniformly in $j$, by strict concavity and compactness) together with the
uniform convergence of Step~3 and convexity of
$\bet\mapsto Q_n(\bet;\lambda,b)$ gives
$\max_{j,b}\|\hat\beta_b(\lambda_j)-\beta(\lambda_j)\|
=O_p(\delta_n)$ with $\delta_n=\sqrt{\log(KBn)/n}$; uniqueness of the
argmax follows from strict concavity of
$\bet\mapsto Q_n(\bet;\lambda,b)$ (its Hessian is
$-\frac1n\sum_i e^{\bet^\top W_i^{(b)}}W_i^{(b)}W_i^{(b)\top}$, positive
semidefinite, and positive definite with probability tending to one by
Assumption~\ref{app:ass:pseudo} and the law of large numbers). Finally,
averaging over $b$ and using convexity,
\[
  \max_{j}\|\bar\beta(\lambda_j)-\beta(\lambda_j)\|
  \le \max_{j,b}\|\hat\beta_b(\lambda_j)-\beta(\lambda_j)\|
  =O_p\!\Bigl(\sqrt{\tfrac{\log(KBn)}{n}}\Bigr)=O_p(n^{-1/2}\,\mathrm{polylog}),
\]
which is $O_p(n^{-1/2})$ up to logarithms; in particular it is
$o_p(1)$, and the refined $O_p(n^{-1/2})$ statement follows from the
linearization in the proof of Theorem~\ref{app:thm:an}. This proves the
lemma.
\end{proof}

\begin{lemma}[Extrapolation approximation error]\label{app:lem:approx}
Grant Assumptions~\ref{app:ass:model}--\ref{app:ass:simgrid}. There exist
constants $C,\rho>1$, depending only on the law of $(X_i,Y_i)$,
$\bSig_\epsilon$, $\bar\lambda$ and the grid, such that for all $d\ge 1$,
\[
  \inf_{\deg(q)\le d}\ \sup_{\lambda\in[-1,\bar\lambda]}
  \bigl\|q(\lambda)-\beta(\lambda)\bigr\|
  \;\le\; C\rho^{-d}.
\]
Consequently, under Assumption~\ref{app:ass:degree} (which forces the selected
degree $\hat d_n\to\infty$), the best degree-$\hat d_n$ approximation
error is $o_p(n^{-1/2})$ whenever $d_{\max}(n)\ge c\log n$ for large $c$.
\end{lemma}

\begin{proof}
By Lemma~\ref{app:lem:curve}(ii), $\lambda\mapsto\beta(\lambda)$ is
real-analytic on $(-1-\eta,\bar\lambda+1)$. Hence each coordinate
$\beta_k(\cdot)$ extends to a holomorphic function on the complex
neighbourhood $D=\{z:|z-\lambda|<\eta \text{ for some }
\lambda\in[-1,\bar\lambda]\}$, in particular on the Bernstein ellipse
$E_r$ with foci $\pm 1$ scaled to the interval $[-1,\bar\lambda]$ and
radius $r>1$. By the standard Chebyshev approximation theorem for
analytic functions (Trefethen, 2013, Theorem~7.1; Rivlin, 1969, Ch.~3), the degree-$d$
polynomial best uniform approximation error on $[-1,\bar\lambda]$
satisfies
$\inf_{\deg q\le d}\sup_{\lambda}|q(\lambda)-\beta_k(\lambda)|
\le M_k r^{-d}$, where
$M_k=\sup_{z\in E_r}|\beta_k(z)|<\infty$. Taking
$\rho=\min(r,1)$-adjusted constants and
$C=\sum_k M_k\vee\ldots$ (map $[-1,\bar\lambda]$ affinely to $[-1,1]$
first; the affine map changes only the constants) proves the first claim.
For the second, $\hat d_n\to\infty$ in probability and
$\hat d_n\le d_{\max}(n)$ give
$\rho^{-\hat d_n}\le\rho^{-c\log n}=n^{-c\log\rho}=o_p(n^{-1/2})$ for
$c$ large enough. (The intersection-of-events argument with
$\Prob(\hat d_n\ge c\log n)\to 1$ completes it.)
\end{proof}

\begin{lemma}[CV-selected degree grows, and selection is stable]\label{app:lem:cv}
Suppose the CV criterion is $\widehat R_n(\lambda_{\max},d)$ computed by
$k$-fold cross-validation over the finite hyperparameter grid
$\mathcal G_n=\Lambda_n^{\max}\times\{1,\dots,d_{\max}(n)\}$, and that:
(i)~$\sup_{g\in\mathcal G_n}|\widehat R_n(g)-R_n(g)|=o_p(1)$, where
$R_n(g)$ is the population CV risk (expected held-out squared prediction
error of the \asimex{} fit with hyperparameters $g$ on a fresh sample);
(ii)~$R_n(\lambda_{\max},d)\ge \rho^{-2d}(1-o(1))$ uniformly over
$d\le d_0$ for any fixed $d_0$ (large-$n$ extrapolation bias dominates
the CV risk at small degrees), and $R_n$ is minimized over $d$ at some
$\bar d_n$ with $\bar d_n\to\infty$ as $K\to\infty$.
Then the CV-selected degree satisfies
$\hat d_n\to\infty$ in probability. If moreover the grid
$\mathcal G_n$ contains a fixed pair $(\lambda^*,d^*)$ and
$R_n(\lambda^*,d^*)=\inf_g R_n(g)+o(n^{-1})$ with a strict population
minimizer, then Assumption~\ref{app:ass:cv} holds.
\end{lemma}

\begin{proof}
Immediate from the triangle inequality:
$|\widehat R_n(g)-R_n(g)|=o_p(1)$ uniformly, so
$\Prob\bigl(\argmin_g\widehat R_n(g)\in\{g:R_n(g)\le\inf_gR_n(g)+\varepsilon\}\bigr)
\to 1$ for every $\varepsilon>0$; combine with (ii) to force
$\hat d_n\to\infty$. The stability claim follows because the minimizer of
$R_n$ over the finite grid is unique for large $n$ and separated from the
rest of the grid by more than $2\sup_g|\widehat R_n(g)-R_n(g)|=o_p(1)$.
Property (i) is a standard cross-validation consistency result
(uniform over a finite grid it follows from the $O_p(n^{-1})$-rate
convergence of held-out losses combined with the union bound); we omit the
routine details. This proves the lemma.
\end{proof}

\begin{remark}\label{app:rem:cvreconcile}
Assumptions~\ref{app:ass:degree} (used for consistency: degree grows) and
\ref{app:ass:cv} (used for the CLT: degree eventually constant) describe two
complementary regimes. The reconciliation is carried out in full in
Section~\ref{app:sec:growing}, where a growing-degree central limit theorem
(Theorem~\ref{app:thm:grow}) is proved under an explicit window condition on
$c\log n$-type degree growth: the extrapolation bias
$\rho^{-d_n}=o(n^{-1/2})$ requires $c\log\rho>1/2$, while control of the
extrapolation variance requires $2c\log\kappa<1$, where
$\kappa=\kappa(\bar\lambda)>1$ is the extrapolation-conditioning constant
of Lemma~\ref{app:lem:extrap}; a valid $c$ exists precisely when the
analyticity margin satisfies $\rho>\kappa$. In the present section we
state the fixed-degree CLT, which is the cleaner result and the one
relevant when the CV degree stabilizes.
\end{remark}

% ===========================================================================
\subsubsection{Theorem 1: Consistency}\label{app:sec:consistency}

\begin{theorem}[Consistency of \asimex{}]\label{app:thm:cons}
Grant Assumptions~\ref{app:ass:model}--\ref{app:ass:degree}. Then
\[
  \hat\beta_{\asimex}\;\xrightarrow{p}\;\bet_0 .
\]
If $d_{\max}(n)\ge c\log n$ eventually, then
$\|\hat\beta_{\asimex}-\bet_0\|=O_p(\rho^{-\hat d_n})
+O_p\bigl(\sqrt{d_{\max}(n)/K(n)}\,n^{-1/2}\bigr)
+O_p(n^{-1/2})$.
\end{theorem}

\begin{proof}
Write $\theta^*_n$ for the population least-squares coefficient
\[
  \theta^*_n=\argmin_{\theta\in\R^{(d+1)p}}
  \frac1K\sum_{j=0}^{K-1}
  \bigl\|\beta(\lambda_j)-\bphi(\lambda_j)^\top\theta\bigr\|^2,
  \qquad p^*(\lambda)=\bphi(\lambda)^\top\theta^*_n,
\]
where we suppress the dependence of $\bphi$ on $\hat d_n$. By
Lemma~\ref{app:lem:unif},
\begin{equation}\label{app:eq:unifbar}
  \Delta_n:=\max_{j}\|\bar\beta(\lambda_j)-\beta(\lambda_j)\|
  =O_p(n^{-1/2}).
\end{equation}

\emph{Step 1: LS coefficient convergence.} The normal equations give
$\hat\theta-\theta^*_n=(\bPhi^\top\bPhi)^{-1}\bPhi^\top
\bar{\bm e}+\bm 0$, where $\bPhi$ is the $K\times(d+1)$ Vandermonde-style
matrix with $j$-th row $\bphi(\lambda_j)^\top$ (same across the $p$
coordinate blocks) and
$\bar{\bm e}=(\bar{\bm e}_0,\dots,\bar{\bm e}_{K-1})$ with
$\bar{\bm e}_j=\bar\beta(\lambda_j)-p^*(\lambda_j)\in\R^p$. Hence
\[
  \|\hat\theta-\theta^*_n\|
  \le\bigl\|(\bPhi^\top\bPhi)^{-1}\bPhi^\top\bigr\|\,
    \|\bar{\bm e}\|_{\mathrm{vec}}
  \;\le\;\frac{1}{\sigma_{\min}(\bPhi)}\,
    \sqrt{K}\,\Delta_n ,
\]
because $\|\bar{\bm e}\|_{\mathrm{vec}}^2=\sum_j\|\bar\beta(\lambda_j)
-p^*(\lambda_j)\|^2\le 2\sum_j\|\bar\beta(\lambda_j)-\beta(\lambda_j)\|^2
+2\sum_j\|\beta(\lambda_j)-p^*(\lambda_j)\|^2$, the second sum being at
most $K$ times the best-approximation error squared. The Gram matrix
$\frac1K\bPhi^\top\bPhi$ converges to the moment matrix of the
degree-$d_{\max}$ polynomial basis under the uniform measure on
$[0,\bar\lambda]$; by classical Vandermonde bounds its minimum eigenvalue
is $\gtrsim d_{\max}^{-C_0}$ for a universal $C_0$ on a grid with
$K\ge 4d_{\max}$ (Assumption~\ref{app:ass:degree}), so
$\sigma_{\min}(\bPhi)\ge c\sqrt{K}\,d_{\max}^{-C_0}$ and
\begin{equation}\label{app:eq:thetahat}
  \|\hat\theta-\theta^*_n\|
  \le C\,d_{\max}^{C_0}\,\Bigl(\Delta_n
     +\sqrt{K}\,\rho^{-\hat d_n}\Bigr).
\end{equation}

\emph{Step 2: extrapolation decomposition.} At $\lambda=-1$,
\[
  \hat\beta_{\asimex}-\bet_0
  =\bphi(-1)^\top(\hat\theta-\theta^*_n)
  +\bigl(p^*(-1)-\beta(-1)\bigr),
  \qquad \beta(-1)=\bet_0,
\]
the last identity by Lemma~\ref{app:lem:curve}(iii). Now
$\|\bphi(-1)\|=\sqrt{\sum_{k=0}^{\hat d_n}1}=\sqrt{\hat d_n+1}$ (since
$(-1)^k=\pm 1$), so by \eqref{app:eq:thetahat},
\[
  \|\bphi(-1)^\top(\hat\theta-\theta^*_n)\|
  \le C\sqrt{\hat d_n}\,d_{\max}^{C_0}\Delta_n
    + C\sqrt{\hat d_n\,K}\,d_{\max}^{C_0}\rho^{-\hat d_n},
\]
and the second term is bounded by the best-approximation error
$C\rho^{-\hat d_n}$ when we note that the least-squares extrapolant at
$-1$ of the \emph{population} curve, $p^*(-1)$, already satisfies
$\|p^*(-1)-\beta(-1)\|\le C\rho^{-\hat d_n}$ by
Lemma~\ref{app:lem:approx}. Collecting terms and using
Lemma~\ref{app:lem:approx} with Assumption~\ref{app:ass:degree}
($\rho^{-\hat d_n}=o_p(n^{-1/2})$ when $d_{\max}\ge c\log n$;
$\Delta_n=O_p(n^{-1/2})$; $\sqrt{\hat d_n}\,d_{\max}^{C_0}$ is polynomial
in $d_{\max}$ and $d_{\max}^3/K\to 0$ with $d_{\max}^2\log n=o(n)$
so that the polynomial factor is absorbed by $n^{-1/2}$):
\[
  \|\hat\beta_{\asimex}-\bet_0\|
  =O_p\!\Bigl(\sqrt{d_{\max}(n)}\,d_{\max}(n)^{C_0}n^{-1/2}\Bigr)
   +O_p\bigl(\rho^{-\hat d_n}\bigr)
  \;=\;o_p(1).
\]
This proves consistency; the displayed rate in the theorem statement
follows from the same inequalities. \qed
\end{proof}

% ===========================================================================
\subsubsection{Theorem 2: Asymptotic normality and the exact variance
decomposition}\label{app:sec:an}

For the CLT we work in the fixed-hyperparameter regime of
Assumption~\ref{app:ass:cv}: $(\hat\lambda_{\max},\hat d)=(\lambda^*,d^*)$
with probability tending to one. To lighten notation write
$d=d^*$, $\lambda_{\max}=\lambda^*$, grid
$\Lambda=\{\lambda_0,\dots,\lambda_{K-1}\}$ (fixed $K$), and
$\bphi(\lambda)=(1,\lambda,\dots,\lambda^d)^\top$.
Define, for $j=0,\dots,K-1$,
\[
  \bm m_{ib}(\lambda_j)
  =\Bigl[Y_i-\exp\bigl(\beta(\lambda_j)^\top W_i^{(b)}(\lambda_j)\bigr)\Bigr]
   W_i^{(b)}(\lambda_j)\in\R^p,
\]
the per-replicate influence score, and
$\bm m_i(\lambda_j)=\E\bigl[\bm m_{ib}(\lambda_j)\bigr]$ (the expectation
is over $(Y_i,X_i,\beps_i,U_{ib})$ jointly and does not depend on $b$).
Note $\E[\bm m_{ib}(\lambda_j)]=\bm 0$ for every $j$, because
$\beta(\lambda_j)$ is the maximizer of
$\E\bigl[\ell(Y_i,W_i^{(b)}(\lambda_j);\bet)\bigr]$ over the joint law
including $U_{ib}$ (Lemma~\ref{app:lem:curve}). Define the matrices
\begin{align}
  \btV^{\mathrm{data}}
  &=\Cov\Bigl(\E\bigl[\bm m_{ib}(\lambda_j)\mid
        Y_i,X_i,\beps_i\bigr]\Bigr)_{j,k=0}^{K-1}
   \in\R^{Kp\times Kp},\label{app:eq:vdata}\\
  \btV^{\mathrm{sim}}
  &=\Cov\Bigl(\bm m_{ib}(\lambda_j)-\E[\bm m_{ib}(\lambda_j)\mid
        Y_i,X_i,\beps_i]\Bigr)_{j,k}
   \in\R^{Kp\times Kp},\label{app:eq:vsim}
\end{align}
i.e.\ the covariance contributed by the data and by the simulated errors,
respectively, both stacked across the grid. Let
\[
  \btH=\operatorname{blockdiag}\bigl(\bm H(\lambda_0),\dots,
        \bm H(\lambda_{K-1})\bigr)\in\R^{Kp\times Kp},
\]
with $\bm H(\lambda)$ the Hessian from Assumption~\ref{app:ass:pseudo}
evaluated at $(\beta(\lambda),1+\lambda)$, and let
$\bm M=\bI_{p}\otimes[(\bPhi^\top\bPhi)^{-1}\bPhi^\top]$ be the
extrapolation map, where $\bPhi\in\R^{K\times(d+1)}$ has rows
$\bphi(\lambda_j)^\top$. Finally
$\bm v_{-1}:=\bI_p\otimes\bphi(-1)$, so that
$\hat\beta_{\asimex}=(\bI_p\otimes\bphi(-1)^\top)\hat\theta
=\bm v_{-1}^\top\bm M\,\bar\beta_{\Lambda}$
with $\bar\beta_{\Lambda}
=(\bar\beta(\lambda_0)^\top,\dots,\bar\beta(\lambda_{K-1})^\top)^\top$.

\begin{theorem}[Asymptotic normality of \asimex{}]\label{app:thm:an}
Grant Assumptions~\ref{app:ass:model}--\ref{app:ass:pseudo}, \ref{app:ass:cv}, with
$B\to\infty$ and fixed $K$, $\lambda_{\max}$, $d$. Then
\[
  \sqrt n\bigl(\hat\beta_{\asimex}-\bet_0\bigr)
  \;\xrightarrow{d}\;N\bigl(\bm 0,\bSig_{\asimex}\bigr),
\]
with the exact decomposition
\begin{equation}\label{app:eq:decomp}
  \bSig_{\asimex}
  \;=\;\bSig_{\mathrm{data}}\;+\;\tfrac1B\,\bSig_{\mathrm{sim}},
  \qquad
  \bSig_{\mathrm{data}}=\bm v_{-1}^\top\bm M\,
    \btH^{-1}\btV^{\mathrm{data}}\btH^{-1}\bm M^\top\bm v_{-1},
\end{equation}
and $\bSig_{\mathrm{sim}}$ defined analogously with
$\btV^{\mathrm{sim}}$ in place of $\btV^{\mathrm{data}}$. In particular,
if $B\to\infty$ then $\bSig_{\mathrm{sim}}/B=o(1)$ and the simulation
component vanishes; the estimator remains asymptotically normal with
variance $\bSig_{\mathrm{data}}$ for \emph{any} rate $B\to\infty$,
including $B=o(n)$.
\end{theorem}

\begin{remark}[Correction to the draft's decomposition
$\bSig_{\mathrm{param}}+\bSig_{\mathrm{extrap}}$]\label{app:rem:decomp}
The draft asserted
$\bSig_{\asimex}=\bSig_{\mathrm{param}}+\bSig_{\mathrm{extrap}}$ with
$\bSig_{\mathrm{param}}=\bm I_W(\bet_0)^{-1}$, the Fisher information of
the \emph{naive} GLM. This is not correct. The \asimex{} estimator is a
weighted linear functional $\sum_j c_j\bar\beta(\lambda_j)$ of the
pseudo-true estimators at the grid points (the weights
$c_j=\bphi(-1)^\top(\bPhi^\top\bPhi)^{-1}\bphi(\lambda_j)$ satisfy
$\sum_j c_j=1$ because the constant function belongs to the basis, so the
functional is bias-correcting), and its asymptotic variance is the
propagated covariance \eqref{app:eq:decomp} of the \emph{influence functions
at the grid points}, not the sum of the naive information and a separate
``extrapolation'' term. The naive information $\bm I_W^{-1}$ is the
variance of the pseudo-estimator at $\lambda=0$ alone; it enters
$\bSig_{\mathrm{data}}$ only through the (perfectly correlated across the
grid) influence structure, and there is no additive split of the
draft's form. The correct additive split, established in
Theorem~\ref{app:thm:an}, is between \emph{data} and \emph{simulation}
sources of randomness, with the simulation term exactly $O(B^{-1})$.
Practically: the draft's variance estimator should be replaced by the
sandwich form \eqref{app:eq:decomp} with $\btV^{\mathrm{data}}$,
$\btV^{\mathrm{sim}}$ estimated from the influence scores
$\bm m_{ib}$; equivalently, a bootstrap over both data and simulated
errors is consistent for $\bSig_{\asimex}$.
\end{remark}

\begin{proof}[Proof of Theorem~\ref{app:thm:an}]
By Assumption~\ref{app:ass:cv} we may and do work on the event
$E_n=\{(\hat\lambda_{\max},\hat d)=(\lambda^*,d^*)\}$, whose probability
tends to one; all statements below are conditional on $E_n$ without
further mention, and the contribution of $E_n^c$ is negligible because
the estimator is bounded in probability (the estimator is a continuous
functional of empirical means on a compact parameter set, hence
$O_p(1)$; formally, use $\Prob(E_n^c)\to 0$ and Slutsky's lemma).

\emph{Step 1: per-replicate linearization.} Fix $j$ and $b$. The map
$\bet\mapsto Q_n(\bet;\lambda_j,b)$ is concave and its Hessian at
$\beta(\lambda_j)$ is
$-\frac1n\sum_i e^{\beta(\lambda_j)^\top W_i^{(b)}}W_i^{(b)}
W_i^{(b)\top}\to_p -\bm H(\lambda_j)$
by the law of large numbers (the summands are i.i.d.\ with finite second
moment; note $W_i^{(b)}(\lambda_j)\stackrel{d}{=}X_i+\beps_\delta$,
$\delta=1+\lambda_j\in[1,1+\bar\lambda]$, so all moments are uniform in
$(j,b)$). The gradient at $\beta(\lambda_j)$ is
$n^{-1}\sum_i\bm m_{ib}(\lambda_j)$ by definition of
$\bm m_{ib}$. A two-term Taylor expansion of the first-order condition
$\nabla_{\bet} Q_n(\hat\beta_b;\lambda_j,b)=\bm 0$ around
$\beta(\lambda_j)$, using the uniform (in $j,b$) convergence of the
Hessian and the $O_p(n^{-1/2}\sqrt{\log(KBn)})$ consistency of
$\hat\beta_b(\lambda_j)$ from Lemma~\ref{app:lem:unif}, yields
\begin{equation}\label{app:eq:lin}
  \sqrt n\bigl(\hat\beta_b(\lambda_j)-\beta(\lambda_j)\bigr)
  =\bm H(\lambda_j)^{-1}\frac{1}{\sqrt n}\sum_{i=1}^n
   \bm m_{ib}(\lambda_j)+o_p(1),
\end{equation}
with the $o_p(1)$ uniform in $(j,b)$ (the remainder is
$O_p(\|\hat\beta_b-\beta(\lambda_j)\|^2)\cdot n
=O_p(\log(KBn))$ at the $\sqrt n$ scale, negligible under
$K^2\log n=o(n)$ and $B\le n^c$; concretely the third derivative of
$\bet\mapsto q(y,w;\bet)$ is bounded by
$e^{C\|w\|}\|w\|^3$, which has uniformly bounded moments, and the
quadratic remainder term divided by $\sqrt n$ is
$O_p(n^{1/2}\delta_n^2)=o(1)$ for
$\delta_n^2=\log(KBn)/n=o(n^{-1/2})$).

\emph{Step 2: averaging over replicates and the two-source
decomposition.} Averaging \eqref{app:eq:lin} over $b=1,\dots,B$,
\[
  \sqrt n\bigl(\bar\beta(\lambda_j)-\beta(\lambda_j)\bigr)
  =\bm H(\lambda_j)^{-1}
   \underbrace{\frac{1}{\sqrt n}\sum_{i=1}^n
   \bar{\bm m}_i(\lambda_j)}_{=:Z_{n,j}}
   +o_p(1),
   \qquad
   \bar{\bm m}_i=\frac1B\sum_{b=1}^B\bm m_{ib},
\]
with the same uniform remainder. Now decompose
$\bar{\bm m}_i=\E[\bm m_{ib}\mid Y_i,X_i,\beps_i]
+\bigl(\bar{\bm m}_i-\E[\bm m_{ib}\mid Y_i,X_i,\beps_i]\bigr)
=:\bm a_i+\bm r_i$,
where $\E[\bm m_{ib}\mid Y_i,X_i,\beps_i]$ does not depend on $b$ (the
conditional law of $U_{ib}$ given $(Y_i,X_i,\beps_i)$ is the same for
every $b$), so $\E[\bm r_i\mid \text{data}]=0$ and, conditionally on the
data, $\bm r_i$ is an average of $B$ independent mean-zero terms with
$\Var(\bm r_i\mid\text{data})=\frac1B\Var(\bm m_{ib}\mid
Y_i,X_i,\beps_i)$; hence
$\Var(\bm r_i)=\frac1B\btV^{\mathrm{sim}}$ in the sense of
\eqref{app:eq:vsim}, and $\Cov(\bm a_i)=\btV^{\mathrm{data}}$ by
\eqref{app:eq:vdata}.

\emph{Step 3: joint CLT across the grid.} The stacked vectors
$\{(\bm a_i,\bm r_i)\}_{i=1}^n$ are i.i.d.\ with finite $2+\nu$ moments
for some $\nu>0$ (Gaussian $\beps,U$ and $\E\|X\|^4<\infty$ provide
uniformly bounded exponential moments of the influence scores; the score
$\bm m_{ib}$ involves $e^{\beta^\top W_i^{(b)}}W_i^{(b)}$ which is
sub-exponential uniformly in $j$). Since $K$ is fixed, the Cram\'er--Wold
device with an arbitrary fixed vector of weights, plus the Lindeberg
condition (verified from the uniform moment bounds), gives
\[
  \frac{1}{\sqrt n}\sum_{i=1}^n
  \begin{pmatrix}\bm a_i\\ \bm r_i\end{pmatrix}
  \;\xrightarrow{d}\;
  N\!\left(\bm 0,\,
  \begin{pmatrix}\btV^{\mathrm{data}} & \bm 0\\ \bm 0 & \bm 0\end{pmatrix}
  \right)
  \quad\text{and}\quad
  \frac{1}{\sqrt n}\sum_{i=1}^n\bm r_i
  \;=\;O_p\!\Bigl(\sqrt{\tfrac1B}\Bigr)\;=\;o_p(1),
\]
because $\Var\bigl(n^{-1/2}\sum_i\bm r_i\bigr)=\btV^{\mathrm{sim}}/B\to
\bm 0$ as $B\to\infty$ (no rate restriction on $B$ beyond
$B\to\infty$ is needed: this corrects the draft's $B/\sqrt n\to\infty$
condition, which was stated without proof and is unnecessary).
Consequently, by the Cram\'er--Wold device applied to the weighted sum
with weights given by the extrapolation map,
\[
  \frac{1}{\sqrt n}\sum_{i=1}^n\bar{\bm m}_{i,\Lambda}
  \;\xrightarrow{d}\;N\bigl(\bm 0,\btV^{\mathrm{data}}\bigr),
  \qquad
  \bar{\bm m}_{i,\Lambda}=\bigl(\bar{\bm m}_i(\lambda_0)^\top,\dots,
  \bar{\bm m}_i(\lambda_{K-1})^\top\bigr)^\top .
\]

\emph{Step 4: extrapolation is a fixed linear map.} The least-squares
coefficient satisfies the exact identity (normal equations, since the
design is nonrandom once the grid is fixed)
\[
  \hat\theta=\bm M\,\bar\beta_\Lambda,
  \qquad\text{hence}\qquad
  \hat\beta_{\asimex}
  =\bm v_{-1}^\top\bm M\,\bar\beta_\Lambda
  =\sum_{j=0}^{K-1}c_j\,\bar\beta(\lambda_j),
\]
with $c_j=\bphi(-1)^\top(\bPhi^\top\bPhi)^{-1}\bphi(\lambda_j)$
nonrandom and $\sum_jc_j=1$. Combining with Steps 2--3,
\[
  \sqrt n\bigl(\hat\beta_{\asimex}-\bet_0\bigr)
  =\bm v_{-1}^\top\bm M\,\btH^{-1}
   \frac{1}{\sqrt n}\sum_{i=1}^n\bar{\bm m}_{i,\Lambda}
  +\bm v_{-1}^\top\bm M\,\btH^{-1}
   \frac{1}{\sqrt n}\sum_{i=1}^n\bm r_{i,\Lambda}
  +\underbrace{\sqrt n\bigl(p^*(-1)-\beta(-1)\bigr)}
   _{=\sqrt n\,O(\rho^{-d})\,=\,o_p(1)}
  +o_p(1),
\]
where $p^*$ is the population LS fit of the (fixed, analytic) curve
$\beta(\cdot)$ by a fixed-degree polynomial; the residual
$p^*(-1)-\beta(-1)$ is the fixed-degree best-approximation error, a
nonrandom constant $O(\rho^{-d})$ (Lemma~\ref{app:lem:approx}) that does not
grow with $n$. \emph{Important:} with fixed $d$ this bias does not
vanish, so normality centered at $\bet_0$ requires the extrapolation
bias to be negligible; this is exactly why Assumption~\ref{app:ass:cv}
(implemented with a degree grid whose population risk is minimized at a
degree $d^*$ where the residual bias is $o(n^{-1/2})$, cf.\
Remark~\ref{app:rem:cvreconcile}) is imposed. Under Assumption~\ref{app:ass:cv},
$\sqrt n\bigl(p^*(-1)-\beta(-1)\bigr)=o_p(1)$.

The first two terms are independent-sources sums; Slutsky's lemma and the
CLT of Step~3 give
\[
  \sqrt n\bigl(\hat\beta_{\asimex}-\bet_0\bigr)
  \xrightarrow{d} N\Bigl(\bm 0,\;
  \bm v_{-1}^\top\bm M\btH^{-1}
  \bigl(\btV^{\mathrm{data}}+\tfrac1B\btV^{\mathrm{sim}}\bigr)
  \btH^{-1}\bm M^\top\bm v_{-1}\Bigr),
\]
where $\btV^{\mathrm{sim}}/B\to\bm 0$ as $B\to\infty$. This is
\eqref{app:eq:decomp}. \qed
\end{proof}

\begin{corollary}[Variance estimation]\label{app:cor:var}
Under the conditions of Theorem~\ref{app:thm:an}, a consistent estimator of
$\bSig_{\asimex}$ is obtained by replacing
$\bm H(\lambda_j)$ with its sample analogue, and
$\btV^{\mathrm{data}},\btV^{\mathrm{sim}}$ with the empirical covariance
matrices of the vectors
$\{\E[\bm m_{ib}\mid Y_i,X_i,\beps_i]\}_i$ and
$\{\bm m_{ib}-\E[\bm m_{ib}\mid Y_i,X_i,\beps_i]\}_{i,b}$
(the conditional expectation is estimated by the sample average over
$b=1,\dots,B$ for each $i$). Equivalently, the pairs-bootstrap over
observations $i$ (holding the simulated errors fixed and redrawing both
data and simulated errors in each bootstrap world) is consistent for
$\bSig_{\asimex}$.
\end{corollary}

\begin{proof}
Each component converges by the same uniform LLN as in
Lemma~\ref{app:lem:unif} (the influence scores have finite second moments
uniformly in $j,b$; the conditional-mean estimator averages $B$
independent terms, giving error $O_p((nB)^{-1/2})=o_p(n^{-1/2})$). The
bootstrap claim follows from standard M-estimation bootstrap theory
applied to the stacked influence representation of Step~2 of
Theorem~\ref{app:thm:an}. Details are routine and omitted.
\end{proof}

% ===========================================================================
\subsubsection{Theorem 4: Asymptotic normality in the growing-degree regime}
\label{app:sec:growing}

This section reconciles the growing-degree regime required for
consistency (Assumption~\ref{app:ass:degree}) with a central limit theorem
centered at $\bet_0$, closing the gap left by Remark~\ref{app:rem:cvreconcile}
of the earlier draft. Throughout, the grid is equispaced on
$[0,\bar\lambda]$ with $K=K(n)$ points and the extrapolation degree is
$d=d(n)$. Map $[0,\bar\lambda]$ affinely to $[-1,1]$ by
$t(\lambda)=2\lambda/\bar\lambda-1$, so that the extrapolation point
$\lambda=-1$ maps to
\[
  t_0=t(-1)=-1-\tfrac{2}{\bar\lambda}<-1,
  \qquad
  \kappa:=|t_0|+\sqrt{t_0^2-1}
  =1+\tfrac{2}{\bar\lambda}
   +\sqrt{\bigl(1+\tfrac{2}{\bar\lambda}\bigr)^2-1}\;>\;1 .
\]
The constant $\kappa$ is the usual Bernstein-parameter of the point $t_0$
outside $[-1,1]$: by the classical Chebyshev extremal property, any
polynomial $q$ of degree $k$ satisfies
$|q(t_0)|\le\kappa^{k}\sup_{t\in[-1,1]}|q(t)|$
(Rivlin, 1969, Ch.~3). Recall from Lemma~\ref{app:lem:approx} that
$\beta(\cdot)$ is real-analytic on a neighbourhood of $[-1,\bar\lambda]$
and admits degree-$k$ approximation error $O(\rho^{-k})$ on
$[-1,\bar\lambda]$, where $\rho>1$ depends on the analyticity strip of
$\beta$; the same holds for the influence functions introduced below.

\begin{lemma}[Extrapolation operator bounds]\label{app:lem:extrap}
Let $K\ge c_0d^2$ and let $\lambda_0<\cdots<\lambda_{K-1}$ be equispaced
on $[0,\bar\lambda]$. Let $\{\varphi_k\}_{k=0}^{d}$ be the polynomials
orthonormal with respect to the discrete inner product
$\langle f,g\rangle_K=\frac1K\sum_jf(\lambda_j)g(\lambda_j)$, and let
\[
  L_{K,d}(h):=\sum_{j=0}^{K-1}c_jh(\lambda_j),
  \qquad
  c_j=\sum_{k=0}^{d}\varphi_k(t_0)\varphi_k(\lambda_j),
\]
be the linear functional ``least-squares degree-$d$ fit evaluated at
$\lambda=-1$''. Then:
\begin{enumerate}\itemsep2pt
\item[(i)] $\displaystyle\sum_{k=0}^{d}\varphi_k(t_0)^2\le C\,d\,\kappa^{2d}$.
\item[(ii)] $\displaystyle\sum_j|c_j|\le C\,d^{1/2}K\kappa^{d}$ and
      $L_{K,d}(q)=q(-1)$ for every polynomial $q$ of degree $\le d$
      (exactness on the model space).
\item[(iii)] For every $h$ real-analytic on a neighbourhood of
      $[-1,\bar\lambda]$ with analytic continuation bounded by $M$ on the
      Bernstein ellipse $E_\rho$ (in $t$-coordinates),
      \[
        \bigl|L_{K,d}(h)-h(-1)\bigr|\;\le\;C\,M\,d^{3/2}K\,\kappa^{d}\rho^{-d}.
      \]
\end{enumerate}
\end{lemma}

\begin{proof}
(i) The discrete and continuous (uniform measure on $[0,\bar\lambda]$)
inner products agree on polynomials of degree $\le 2d$ up to relative
error $O(d^2/K)$: for $p\ge0$ of degree $\le 2d$,
$\bigl|\frac1K\sum_jp(\lambda_j)-\int p\,d\mu\bigr|
\le (C d^2/K)\int p\,d\mu$,
a classical composite-quadrature estimate combining the
Euler--Maclaurin error on each subinterval with the Markov inequality
$\|p''\|_\infty\le C d^4\|p\|_\infty\le C d^4(d+1)^2\int p\,d\mu$
(Rivlin, 1969, Ch.~3; the last step is the $L^1$--$L^\infty$ Nikolskii
inequality for degree-$2d$ polynomials). Consequently the Gram matrices
of $\{1,t,\dots,t^d\}$ under the two inner products are equivalent with
relative error $O(d^2/K)=O(1/c_0)$, so the discrete Christoffel function
is within a constant factor of the continuous one:
\[
  \sum_{k=0}^{d}\varphi_k(t_0)^2
  \;\le\;C\sum_{k=0}^{d}\bigl(k+\tfrac12\bigr)P_k(t_0)^2
  \;\le\;C\,d\,\kappa^{2d},
\]
where $P_k$ are the Legendre polynomials (orthonormal versions
$\tilde p_k=\sqrt{k+1/2}\,P_k$), using $|P_k(t_0)|\le\kappa^k$ from the
Chebyshev extremal property and $\sum_{k\le d}(k+\tfrac12)\le C d^2$
against the dominant term $\kappa^{2d}$.

(ii) By the reproducing-kernel identity
$c_j=\sum_k\varphi_k(t_0)\varphi_k(\lambda_j)$ and discrete
orthonormality $\sum_j\varphi_k(\lambda_j)\varphi_{k'}(\lambda_j)
=K\delta_{kk'}$,
\[
  \sum_jc_j^2
  =\sum_{k,k'}\varphi_k(t_0)\varphi_{k'}(t_0)
   \textstyle\sum_j\varphi_k(\lambda_j)\varphi_{k'}(\lambda_j)
  =K\sum_k\varphi_k(t_0)^2
  \le C\,d\,K\kappa^{2d},
\]
whence $\sum_j|c_j|\le\sqrt K\,\bigl(\sum_jc_j^2\bigr)^{1/2}
\le C\,d^{1/2}K\kappa^{d}$ by Cauchy--Schwarz. Exactness: $L_{K,d}$ is
the evaluation at $t_0$ of the $L^2$-orthogonal (discrete) projection
onto $\mathcal P_d$, which reproduces every $q\in\mathcal P_d$; and
evaluation is linear, so $L_{K,d}(q)=q(t_0)=q(-1)$ in $\lambda$-terms
(since $q$ as a polynomial in $t$ maps to a polynomial of the same degree
in $\lambda$ and $t(-1)=t_0$).

(iii) Let $q$ be the degree-$d$ polynomial minimizing
$\sup_{[-1,\bar\lambda]}|h-q|\le CM\rho^{-d}$ (Chebyshev approximation of
analytic functions; Lemma~\ref{app:lem:approx}). Then
\[
  |L_{K,d}(h)-h(-1)|
  \le|L_{K,d}(h-q)|+|q(-1)-h(-1)|
  \le\Bigl(\sum_j|c_j|\Bigr)CM\rho^{-d}+CM\rho^{-d},
\]
and (ii) gives the claim.
\end{proof}

Define, for $\lambda\in[0,\bar\lambda]$, the data-influence function
\[
  \bm\psi(\lambda)
  :=\E\bigl[\bm m_{ib}(\lambda)\mid Y_i,X_i,\beps_i\bigr]
  =Y_i\bm w_i
   -\exp\!\Bigl(\beta(\lambda)^\top\bm w_i
   +\tfrac{\lambda\sigma^2}{2}\beta(\lambda)^\top\bSig_0\beta(\lambda)\Bigr)
   \bigl(\bm w_i+\lambda\sigma^2\bSig_0\beta(\lambda)\bigr),
\]
with $\bm w_i=X_i+\beps_i$ (this is the factorized conditional expectation
of $\bm m_{ib}$ from the proof of Lemma~\ref{app:lem:curve}; the exponential
tilt is the $\lambda$-level augmentation with
$U\sim N(0,\sigma^2\bSig_0)$). By Lemma~\ref{app:lem:curve}(ii),
$\lambda\mapsto\beta(\lambda)$ is real-analytic on a neighbourhood of
$[-1,\bar\lambda]$, hence so is $\bm\psi(\lambda)$, and so is the
covariance function
\[
  \bm\Gamma(\lambda,\lambda')
  :=\Cov\bigl(\bm\psi_1(\lambda),\bm\psi_1(\lambda')\bigr),
\]
defined and real-analytic on a neighbourhood of
$[-1,\bar\lambda]^2$. Write $\widetilde{\bm\Gamma}$ for its analytic
continuation and $\widetilde{\bm Z}:=\widetilde{\bm\Gamma}( -1,-1)$;
note $\bm\Gamma(\lambda,\lambda')\to\bm I_X(\bet_0)$ as
$\sigma^2\to0$ uniformly on $[-1,\bar\lambda]^2$ (by the same
$L^2$-collapse as in the proof of Theorem~\ref{app:thm:are}), so
$\widetilde{\bm Z}\to\bm I_X(\bet_0)$.

\begin{theorem}[Growing-degree CLT]\label{app:thm:grow}
Grant Assumptions~\ref{app:ass:model}--\ref{app:ass:simgrid}, with $K=K(n)$
equispaced points and $B=B(n)\to\infty$. Suppose the selected degree
satisfies $\hat d_n/d_n\xrightarrow{p}1$ for a deterministic sequence
\[
  d_n=\lceil c\log n\rceil,
  \qquad K(n)\ge c_0d_n^2,
\]
and that the analyticity margin condition
\begin{equation}\label{app:eq:window}
  \rho>\kappa
\end{equation}
holds with $c$ chosen in the window
$\tfrac{1}{2\log\rho}<c<\tfrac{1}{2\log\kappa}$ (nonempty exactly under
\eqref{app:eq:window}). Then
\[
  \sqrt n\bigl(\hat\beta_{\asimex}-\bet_0\bigr)
  \;\xrightarrow{d}\;N\bigl(\bm 0,\bSig_*\bigr),
  \qquad
  \bSig_*=\bm I_X(\bet_0)^{-1}\,\widetilde{\bm Z}\,
           \bm I_X(\bet_0)^{-1},
\]
with $\widetilde{\bm Z}$ as defined above; the simulation component
contributes $O(B^{-1})$ to $\bSig_*$ and vanishes as $B\to\infty$, and
the extrapolation bias is $o(n^{-1/2})$, so the limit is centered at
$\bet_0$. In particular $\bSig_*\to\bSig_{\mathrm{oracle}}$ as
$\sigma^2\to0$, consistently with Theorem~\ref{app:thm:are}.
\end{theorem}

\begin{remark}
The window condition has a transparent interpretation: $\rho$ measures
how far the analytic continuation of the SIMEX curve extends (bias
side), while $\kappa$ measures how violently degree-$d$ polynomial
extrapolation to a point outside the grid amplifies noise (variance
side); Theorem~\ref{app:thm:grow} is feasible precisely when the former
dominates the latter. If \eqref{app:eq:window} fails, growing-degree
extrapolation is variance-dominated and the fixed-degree CLT
(Theorem~\ref{app:thm:an}) is the appropriate inference tool.
\end{remark}

\begin{proof}[Proof of Theorem~\ref{app:thm:grow}]
\emph{Step 1: bias.} The deterministic extrapolation error is
$|L_{K,d_n}(\beta)-\beta(-1)|\le C d_n^{3/2}K\kappa^{d_n}\rho^{-d_n}$ by
Lemma~\ref{app:lem:extrap}(iii) applied to the (deterministic, analytic,
$M$-bounded) curve $\beta$. In $\sqrt n$ units this is
$\sqrt n\,d_n^{3/2}K\kappa^{d_n}\rho^{-d_n}
= \mathrm{polylog}(n)\;n^{1/2+c\log\kappa-c\log\rho}\to0$
because $c(\log\rho-\log\kappa)>0$. Hence centering at $\bet_0=\beta(-1)$
is valid.

\emph{Step 2: uniform per-grid-point linearization.} Exactly as in the
proof of Theorem~\ref{app:thm:an} (Step~1--2), with the maximal inequality
entropy now carrying $\log(KB)$ and the quadratic Taylor remainder
$n^{1/2}\delta_n^2=O(\log(KBn)/n^{1/2})=o(1)$ (using
$K^2\log n=o(n)$ and $B\le n^{c'}$),
\[
  \sqrt n\bigl(\bar\beta(\lambda_j)-\beta(\lambda_j)\bigr)
  =\bm H(\lambda_j)^{-1}Z_n(\lambda_j)+o_p(1),
  \quad\text{uniformly in }j,
\]
where $Z_n(\lambda)=n^{-1/2}\sum_i\bar{\bm m}_i(\lambda)$ and
$\bar{\bm m}_i(\lambda)=B^{-1}\sum_b\bm m_{ib}(\lambda)$.

\emph{Step 3: empirical-process convergence with analytic limit paths.}
Decompose $\bar{\bm m}_i=\bm\psi_i+\bm r_i$ as in
Theorem~\ref{app:thm:an}: $\bm r_i$ has conditional mean zero and variance
$B^{-1}\btV^{\mathrm{sim}}$, so
$\sup_\lambda\|n^{-1/2}\sum_i\bm r_i(\lambda)\|
=O_p((nB)^{-1/2}\,\mathrm{polylog})=o_p(1)$
(the supremum over the grid is controlled by the same maximal
inequality over the one-dimensional analytic family; the entropy of the
class $\{\bm m_{ib}(\lambda)-\bm\psi(\lambda)\}$ is
$O(\log(1/\delta))$). The class
$\{\bm\psi(\lambda):\lambda\in[0,\bar\lambda]\}$ is a one-dimensional
analytic family with square-integrable envelope; its bracketing entropy
is $O(\log(1/\delta))$, whence (i)~the functional CLT
\[
  Z_n(\cdot)-n^{-1/2}\textstyle\sum_i\bm r_i(\cdot)
  \;\Rightarrow\;Z(\cdot)\quad\text{in }\ell^\infty[0,\bar\lambda],
\]
with $Z$ a centered Gaussian process with covariance $\bm\Gamma$, and
(ii)~the sup-norm rate
$\|Z_n-Z\|_\infty=O_p\bigl(\sqrt{\log n/n}\bigr)$
(van der Vaart and Wellner, 1996, Thms.~2.14.1--2.14.2). The
sample paths of $Z$ are almost surely real-analytic on a neighbourhood
of $[-1,\bar\lambda]$ with analytically continuable coordinates: this
follows from the analyticity of the mean (zero) and covariance
($\bm\Gamma$) and the classical result that Gaussian processes with
analytic covariance in a strip have a.s.\ analytic paths there
(e.g.\ via the Cauchy--Kovalevskaya-type majorization of
$\E\|\partial_\lambda^k Z\|^2\le C^k k!$, which ensures the Taylor
series converges a.s.).

\emph{Step 4: reduction to a single linear functional.} By Step~2,
\[
  \sqrt n\bigl(\hat\beta_{\asimex}-\bet_0\bigr)
  =L_{K,d_n}\bigl(\bm H(\cdot)^{-1}Z_n(\cdot)\bigr)
   +o_p(1)+\underbrace{\sqrt n\,L_{K,d_n}(\beta-\beta(-1))}_{\to0,\ \text{Step 1}},
\]
because $\hat\beta_{\asimex}-\bet_0=L_{K,d_n}(\bar\beta-\beta(-1))$
(exactness of $L_{K,d_n}$ on polynomials, Lemma~\ref{app:lem:extrap}(ii))
and $\bar\beta-\beta(-1)=[\bar\beta-\beta]+[\beta-\beta(-1)]$ with the
first bracket linearized in Step~2 and the second treated in Step~1.

\emph{Step 5: convergence of the functional.} Write
$\bm g_\omega(\lambda)=\bm H(\lambda)^{-1}Z(\lambda,\omega)$ for the
(a.s.\ analytic) coordinate paths. On the event
$E_n=\{$the analytic continuations of all $p$ coordinates of
$\bm g_\omega$ are bounded by $M_n$ on $E_\rho\}$, where
$M_n$ is any sequence with $M_n\to\infty$ slowly
(e.g.\ $M_n=\log\log n$), Lemma~\ref{app:lem:extrap}(iii) gives
\[
  \sup_{\omega\in E_n}\bigl|L_{K,d_n}(\bm g_\omega)
  -\bm H(-1)^{-1}\widetilde Z(-1,\omega)\bigr|
  \le C M_n d_n^{3/2}K\kappa^{d_n}\rho^{-d_n}\longrightarrow 0,
\]
by the same computation as in Step~1; and
$\Prob(E_n)\to1$ by the Borell--TIS inequality applied to the
supremum of the analytic continuation of the Gaussian process over the
ellipse (a.s.\ finite, with Gaussian concentration).
The discretization error
$L_{K,d_n}\bigl(\bm H^{-1}(Z_n-Z)\bigr)$ satisfies
\[
  \Bigl\|L_{K,d_n}\bigl(\bm H^{-1}(Z_n-Z)\bigr)\Bigr\|
  \le\Bigl(\sup_j\|\bm H(\lambda_j)^{-1}\|\Bigr)
     \Bigl(\sum_j|c_j|\Bigr)\|Z_n-Z\|_\infty
  =O\bigl(d_n^{1/2}K\kappa^{d_n}\bigr)\,O_p\!\Bigl(\sqrt{\tfrac{\log n}{n}}\Bigr)
  \to0,
\]
because $d_n^{1/2}K\kappa^{d_n}\le\mathrm{poly}(d_n)\,
n^{c\log\kappa}=o\bigl(n^{1/2}/\sqrt{\log n}\bigr)$ under
$2c\log\kappa<1$. Finally, the linearization remainder of Step~2
propagated through $L_{K,d_n}$ is $o_p(1)$ by the same operator bound.

\emph{Step 6: the limit law.} Combining Steps 4--5,
\[
  \sqrt n\bigl(\hat\beta_{\asimex}-\bet_0\bigr)
  =\bm H(-1)^{-1}\widetilde Z(-1)+o_p(1),
\]
and it remains to identify the law of the right-hand side. For each $n$
the vector $\bm H(-1)^{-1}L_{K,d_n}\bigl(\bm H^{-1}Z_n\bigr)$ is a
centered sample average of i.i.d.\ bounded-$(2+\nu)$-moment vectors
divided by $\sqrt n$, with per-observation contributions bounded by
$O\bigl(d_n^{1/2}K\kappa^{d_n}\bigr)/\sqrt n=o(1)$ times a fixed
sub-exponential variable; the Lindeberg condition therefore holds, and
its variance converges by the double functional convergence
\[
  \Cov\Bigl(L_{K,d_n}\bigl(\bm H^{-1}Z_n\bigr)\Bigr)
  =\sum_{j,k}c_jc_k\,\bm H_j^{-1}\bm\Gamma(\lambda_j,\lambda_k)\bm H_k^{-1}
  \longrightarrow
  \bm H(-1)^{-1}\widetilde{\bm Z}\,\bm H(-1)^{-1},
\]
applying Lemma~\ref{app:lem:extrap}(iii) twice to the real-analytic
function $(\lambda,\lambda')\mapsto
\bm H(\lambda)^{-1}\bm\Gamma(\lambda,\lambda')\bm H(\lambda')^{-1}$
(exactness on polynomials is not needed here, only (iii), which applies
to analytic functions of two variables by the same Chebyshev argument on
the product ellipse). By the Cram\'er--Wold device the limit is
$N(\bm 0,\bSig_*)$ with $\bSig_*=\bm I_X^{-1}\widetilde{\bm Z}\bm
I_X^{-1}$, using $\bm H(-1)=\bm H(\bet_0,0)=\bm I_X(\bet_0)$. The
simulation component adds $O(B^{-1})$ to the variance, as in
Theorem~\ref{app:thm:an}. \qed
\end{proof}

% ===========================================================================
\subsubsection{Theorem 3: Asymptotic relative efficiency}\label{app:sec:are}

Let $\bSig_{\mathrm{oracle}}=\bm I_X(\bet_0)^{-1}$ with
$\bm I_X(\bet_0)=\E\bigl[e^{X_i^\top\bet_0}X_iX_i^\top\bigr]$, and define
the generalized asymptotic relative efficiency
\[
  \mathrm{ARE}(\asimex\text{ vs.\ Oracle})
  =\frac{\det\bigl(\bSig_{\mathrm{oracle}}\bigr)}
        {\det\bigl(\bSig_{\asimex}\bigr)} .
\]
Three information matrices must be distinguished. (a)~The oracle
information $\bm I_X(\bet_0)$, for the model of $(Y_i,X_i)$. (b)~The
information $\bm I(t)$ of the \emph{observed-data} model
$(Y_i,\widetilde W_t)$, $\widetilde W_t=X_i+\bm e_t$ with
$\bm e_t\sim N(\bm 0,t\bSig_0)$, $t\ge0$ (so $t=\sigma^2$ in the
application); this is the information that bounds what \emph{any}
estimator based on noisy covariates can achieve. (c)~The information
$\E[e^{\widetilde W^\top\bet_0}\widetilde W\widetilde W^\top]$ of the
\emph{naive working model}, which is what the draft denoted
$\bm I_W$ and used as $\bSig_{\mathrm{param}}$; it is a third object,
the curvature of a misspecified likelihood, and plays no role in the
correct theory (see Remark~\ref{app:rem:decomp}).

\begin{lemma}[Information monotonicity under additive noise]\label{app:lem:info}
Grant Assumptions~\ref{app:ass:model} with $\bSig_0$ fixed. Then:
\begin{enumerate}\itemsep2pt
\item[(i)] For $0\le t_1\le t_2$, $\bm I(t_2)\preceq\bm I(t_1)$ in the
      Loewner order; in particular $\bm I(t)\preceq\bm I(0)=\bm I_X(\bet_0)$
      for all $t\ge0$.
\item[(ii)] The inequality is strict for $t>0$: $\bm I(t)\prec\bm I_X(\bet_0)$
      provided $X_i$ is not a.s.\ constant.
\end{enumerate}
Consequently $\bSig_{\mathrm{oracle}}=\bm I_X(\bet_0)^{-1}
\preceq\bm I(t)^{-1}$ for all $t\ge0$: the oracle covariance is the
smallest, and by the H\'ajek convolution theorem
$\bm I(t)^{-1}$ is the best asymptotic covariance achievable by any
regular estimator based on $(Y_i,\widetilde W_t)$ alone.
\end{lemma}

\begin{proof}
For $t_2>t_1$ write $\widetilde W_{t_2}
=\widetilde W_{t_1}+\bm\eta$ with
$\bm\eta\sim N(\bm 0,(t_2-t_1)\bSig_0)$ independent of
$(Y_i,\widetilde W_{t_1})$; the model of $(Y,\widetilde W_{t_2})$ is
thus the image of the model of $(Y,\widetilde W_{t_1})$ under a
Markov kernel that does not depend on $\bet$. Standard score
transformation under parameter-independent smoothing: if $Z\sim P_{\bet}$
and $Z'=Z+\bm\eta$ with $\bm\eta\perp\bet$, then the score of $Z'$
equals the conditional expectation of the score of $Z$ given $Z'$,
\[
  \bm s_{Z'}(\bet)=\E\bigl[\bm s_Z(\bet)\bigm|Z'\bigr],
\]
justified by differentiating
$p_{Z'}(z')=\int k(z'\leftarrow z)\,p_Z(z)\,dz$ under the integral:
domination holds because the Gaussian kernel supplies exponential
moments and the score
$\bm s_Z(\bet)=[Y-e^{\bet^\top X}]X$ (for $Z=(Y,X)$) has uniformly
bounded exponential moments. Then, by the law of total variance,
\[
  \bm I(t_2)=\Var\bigl(\E[\bm s_Z\mid Z']\bigr)
  \;\preceq\;\Var(\bm s_Z)=\bm I(t_1),
\]
with equality iff $\bm s_Z$ is $Z'$-measurable. For $t_1=0$,
$\bm s_Z=[Y-e^{\bet_0^\top X}]X$ with $Z=(Y,X)$; if equality held for
some $t_2>0$, $\bm s_Z$ would be measurable with respect to
$(Y,X+\bm e_{t_2})$, which is impossible unless $X$ is a.s.\ constant
(the conditional law of $X$ given $(Y,X+\bm e_t)$ is non-degenerate for
$t>0$ on a set of positive probability when $X$ is not a.s.\ constant,
and $\bm s_Z$ genuinely depends on $X$). This proves (i)--(ii). The
final display follows by inverting the Loewner order (for positive
definite matrices, $\bm 0\prec\bm A\preceq\bm B$ implies
$\bm B^{-1}\preceq\bm A^{-1}$) and invoking the H\'ajek convolution
theorem for the local limit of regular estimators in the model of
$(Y,\widetilde W_t)$.
\end{proof}

\begin{theorem}[Efficiency relative to the oracle]\label{app:thm:are}
Grant the conditions of Theorem~\ref{app:thm:an} with $B\to\infty$, and
suppose $\bSig_\epsilon=\sigma^2\bSig_0$ with $\bSig_0$ fixed and
$\sigma^2\downarrow0$ (with all other quantities, including the grid and
the law of $X_i$, held fixed). Write
$\bSig_{\asimex}(\sigma^2)$ for the asymptotic covariance of
Theorem~\ref{app:thm:an} at measurement-error scale $\sigma^2$. Then:
\begin{enumerate}\itemsep2pt
\item[(i)] \textbf{(First-order efficiency.)}
      $\bSig_{\asimex}(\sigma^2)\to\bSig_{\mathrm{oracle}}$
      as $\sigma^2\to0$, for every fixed grid and degree: the limit is
      exact rather than merely $O(\sigma^2)$-close.
\item[(ii)] \textbf{(Explicit second-order term.)}
      $\bSig_{\asimex}(\sigma^2)
      =\bSig_{\mathrm{oracle}}+\sigma^2\,\dot{\bm C}_\beta+O(\sigma^4)$,
      where $\dot{\bm C}_\beta$ is the \emph{symmetric}, explicitly
      computable matrix
      \begin{equation}\label{app:eq:cdot}
        \dot{\bm C}_\beta
        =\bm v_{-1}^\top\bm M\,\bm H_0^{-1}
        \Bigl(\dot{\btV}-\dot{\btH}\bm H_0^{-1}\bm V_0
              -\bm V_0\bm H_0^{-1}\dot{\btH}\Bigr)
        \bm H_0^{-1}\bm M^\top\bm v_{-1},
      \end{equation}
      with $\bm H_0=\bI_K\otimes\bm I_X(\bet_0)$,
      $\bm V_0=(\bm 1_K\bm 1_K^\top)\otimes\bm I_X(\bet_0)$, and
      \begin{align}
        \dot{\btH}_j
        &=\tfrac{\delta_j}{2}\bigl(\bet_0^\top\bSig_0\bet_0\bigr)\bm I_X
          +\delta_j\,\E\bigl[e^{\bet_0^\top X}
             (\tilde{\bm b}_{\bSig_0}^\top X)XX^\top\bigr]
          +\bSig_0\bet_0\,\bm g_0^\top+\bm g_0\,\bet_0^\top\bSig_0,
          \label{app:eq:hdot}\\
        \dot{\btV}_{jk}
        &=\Cov\bigl(\dot{\bm\psi}_j,\bm\psi^{(0)}\bigr)
         +\Cov\bigl(\bm\psi^{(0)},\dot{\bm\psi}_k\bigr),
        \qquad
        \dot{\bm\psi}_j
        =-e^{\bet_0^\top X}
         \Bigl[\bigl(\delta_j\tilde{\bm b}_{\bSig_0}^\top X
                +\tfrac{\lambda_j}{2}\bet_0^\top\bSig_0\bet_0\bigr)X
               +\lambda_j\bSig_0\bet_0\Bigr],
        \label{app:eq:vdot}
      \end{align}
      where $\delta_j=1+\lambda_j$,
      $\bm g_0=\E[e^{\bet_0^\top X}X]$, and
      $\bm\psi^{(0)}=YX-e^{\bet_0^\top X}X$ is the oracle influence
      function. Symmetry of $\dot{\btH}_j$ and $\dot{\btV}$ is immediate
      from \eqref{app:eq:hdot}--\eqref{app:eq:vdot}, and
      \eqref{app:eq:cdot} is a congruence-like transform of symmetric
      matrices, hence $\dot{\bm C}_\beta$ is symmetric.
      In general, however, $\dot{\bm C}_\beta$ is
      \emph{indefinite}: neither
      $\dot{\bm C}_\beta\succeq\bm 0$ nor
      $\dot{\bm C}_\beta\preceq\bm 0$ holds universally, and the
      sign of
      $\operatorname{tr}\bigl(\bSig_{\mathrm{oracle}}^{-1}
      \dot{\bm C}_\beta\bigr)$ depends on the design, the grid, and the
      degree (see Remark~\ref{app:rem:psd}).
\item[(iii)] \textbf{(Determinant expansion and correct efficiency
      statement.)}
      \[
        \mathrm{ARE}=1+\sigma^2\kappa_\beta+O(\sigma^4),
        \qquad
        \kappa_\beta:=-\operatorname{tr}\!\bigl(
          \bSig_{\mathrm{oracle}}^{-1}\dot{\bm C}_\beta\bigr),
      \]
      with no universal sign on $\kappa_\beta$. The universal
      statements are: (a)~$\mathrm{ARE}\to1$ as $\sigma^2\to0$
      (part~(i)); (b)~by Lemma~\ref{app:lem:info},
      $\bSig_{\mathrm{oracle}}\preceq\bm I(\sigma^2)^{-1}$, so no
      regular estimator based on the noisy covariates alone can improve
      on the oracle to first order; and (c)~\asimex{} attains this
      frontier to first order, its excess over
      $\bSig_{\mathrm{oracle}}$ being second-order
      ($O(\sigma^2)$) with a bias--variance trade-off reflected in the
      sign of $\kappa_\beta$.
\end{enumerate}
\end{theorem}

\begin{remark}[Why the draft's positive semidefiniteness claim is
false]\label{app:rem:psd}
The draft asserted $\dot{\bm C}_\beta\succeq\bm 0$, hence
$\mathrm{ARE}\le1+O(\sigma^4)$. This is not correct, for a structural
reason: measurement error attenuates the signal but can also
\emph{de-noise} the estimating directions. In the homoskedastic linear
model $Y=X\bet_0+e$, $W=X+\epsilon$, the naive slope has asymptotic variance
$\sigma_e^2/\E[W^2]<\sigma_e^2/\E[X^2]$, the oracle variance: the
pseudo-estimators' variances \emph{decrease} in $\sigma^2$, and the
resulting negative contribution enters $\dot{\btV}$ through
\eqref{app:eq:vdot}. Since \asimex{} is a fixed linear combination
($\sum_jc_j=1$) of these pseudo-estimators, its second-order variance
term inherits this sign. The Poisson log-link exhibits the same
phenomenon for suitable designs. What survives universally is
part~(iii): the oracle marks the first-order frontier
(Lemma~\ref{app:lem:info}), \asimex{} reaches it, and the second-order gap
splits into the extrapolation bias
$\sigma^2\bigl(\sum_jc_j\delta_j\bigr)\tilde{\bm b}_{\bSig_0}$
(deterministic, $O(\sigma^2)$, vanishing for grids with
$\sum_jc_j\delta_j=0$) plus the symmetric variance correction
\eqref{app:eq:cdot} of indefinite sign.
\end{remark}

\begin{corollary}[Explicit scalar constants]\label{app:cor:scalar}
For $p=1$, the first-order naive bias is, by \eqref{app:eq:bias1},
\[
  \tilde b=-\frac{\tfrac12\beta_0^{2}\,\mu_1(\beta_0)
        +\beta_0\,\mu_0(\beta_0)}{\mu_2(\beta_0)}\cdot\sigma^2\text{-scale},
  \qquad \mu_r(\beta)=\E[e^{\beta X}X^r],
\]
so for $X\sim N(0,1)$,
$\tilde b=-\sigma^2\beta_0\bigl(1+\beta_0^2/2\bigr)/(1+\beta_0^2)$
(attenuation). The second-order constant is
\[
  \kappa_\beta
  =-\frac{1}{\mu_2(\beta_0)}
   \sum_{j,k}c_jc_k\Bigl[
     \dot\Gamma_{jk}
     -\dot H_j\,\Gamma_{jk}^{(0)}\cdot\frac{1}{\mu_2(\beta_0)}
     -\Gamma_{jk}^{(0)}\dot H_k\cdot\frac{1}{\mu_2(\beta_0)}
   \Bigr],
\]
with $\Gamma_{jk}^{(0)}=\Var(\bm\psi^{(0)})=\mu_2(\beta_0)$,
$\dot H_j$ from \eqref{app:eq:hdot} (scalar), and $\dot\Gamma_{jk}$ from
\eqref{app:eq:vdot} (scalar); every entry is an explicit functional of
$\{\mu_r(\beta_0)\}_{r\le4}$ and the grid weights $\{c_j\}$.
\end{corollary}

\begin{proof}[Proof of Theorem~\ref{app:thm:are}]
\emph{Step 1: the $\sigma^2\to0$ collapse of the influence structure.}
As $\sigma^2\downarrow0$, for every grid point $\lambda_j$,
$\beta(\lambda_j)\to\beta(0)=\bet_0$
(Lemma~\ref{app:lem:curve}(ii)--(iii) with $\delta_j=1+\lambda_j\to1$, and
$\beta(\delta)\to\beta(0)$ by continuity; note $\beta(0)=\bet_0$
corresponds to the oracle). The influence score at $\lambda_j$,
\[
  \bm m_{ib}(\lambda_j)
  =\bigl[Y_i-e^{\beta(\lambda_j)^\top W_i^{(b)}}\bigr]
   W_i^{(b)},\qquad
  W_i^{(b)}=X_i+\beps_i+\sqrt{\lambda_j}U_{ib},
\]
converges in $L^2$ to $\bm m_i(0)=
\bigl[Y_i-e^{\bet_0^\top X_i}\bigr]X_i$, because
$W_i^{(b)}\to X_i$ in $L^4$ (Gaussian tails),
$e^{\beta(\lambda_j)^\top W_i^{(b)}}\to e^{\bet_0^\top X_i}$ in $L^2$
(locally uniform convergence plus dominated convergence with the
sub-exponential dominant $e^{C\|W_i^{(b)}\|}$, whose moments are bounded
uniformly in $\sigma^2\le\bar\sigma^2$), and the products converge by
Cauchy--Schwarz. Hence the stacked covariance converges to
$\bm V_0=(\bm 1_K\bm 1_K^\top)\otimes\bm I_X(\bet_0)$ (perfect
cross-grid correlation: all $K$ influence vectors collapse to the same
oracle influence) and $\btH\to\bI_K\otimes\bm I_X(\bet_0)=\bm H_0$.

\emph{Step 2: the extrapolation functional becomes the identity.} The
extrapolated estimator is
$\hat\beta_{\asimex}=\sum_j c_j\bar\beta(\lambda_j)$ with
$\sum_j c_j=1$. Substituting the limit covariance,
\[
  \bSig_{\asimex}\to
  \sum_{j,k}c_jc_k\cdot\bm I_X(\bet_0)^{-1}
  \bm I_X(\bet_0)\bm I_X(\bet_0)^{-1}
  =\Bigl(\sum_jc_j\Bigr)^2\bm I_X(\bet_0)^{-1}
  =\bSig_{\mathrm{oracle}},
\]
since $\sum_jc_j=1$ (the constant function lies in the span of the
basis). Equivalently, in matrix form:
$\btH_0^{-1}\bm V_0\btH_0^{-1}
=(\bm 1\bm 1^\top)\otimes\bm I_X^{-1}$ and
$\bm v_{-1}^\top\bm M\bigl((\bm 1\bm 1^\top)\otimes\bm I_X^{-1}\bigr)
\bm M^\top\bm v_{-1}=\bigl(\sum_jc_j\bigr)^2\bm I_X^{-1}
=\bm I_X^{-1}$. This proves (i).

\emph{Step 3: first-order expansion of the ingredients.} All
$\sigma^2$-dependence enters through (a) the pseudo-true curve
$\beta(\lambda_j;\sigma^2)=\bet_0+\delta_j\sigma^2
\tilde{\bm b}_{\bSig_0}+O(\sigma^4)$ (Lemma~\ref{app:lem:curve}(iii)),
with
\[
  \tilde{\bm b}_{\bSig_0}
  =-\bm I_X(\bet_0)^{-1}\Bigl[
    \tfrac12\bigl(\bet_0^\top\bSig_0\bet_0\bigr)\,
    \E\bigl[e^{\bet_0^\top X}X\bigr]
    +\bSig_0\bet_0\,\E\bigl[e^{\bet_0^\top X}\bigr]\Bigr],
\]
(b) the Gaussian tilt factor
$e^{(\lambda_j\sigma^2/2)\beta^\top\bSig_0\beta}$ in the conditional
mean of the influence scores, and (c) the tilt of the design moment
$\E[e^{\beta^\top(X+\beps_\delta)}(X+\beps_\delta)(X+\beps_\delta)^\top]$.
For (c), the factorization
\[
  \bm H_j(\sigma^2)
  =e^{\frac{\delta_j\sigma^2}{2}\beta_j^\top\bSig_0\beta_j}
   \Bigl[\bm H(\beta_j,0)
    +\delta_j\sigma^2\bigl(\bSig_0\beta_j\bm g_j^\top
      +\bm g_j\beta_j^\top\bSig_0\bigr)\Bigr],
\]
(where $\bm H(\beta,0)=\E[e^{\beta^\top X}XX^\top]$,
$\bm g_j=\E[e^{\beta_j^\top X}X]$, and we used
$\E[e^{\beta^\top\beps_\delta}\beps_\delta]
=\delta\sigma^2\bSig_0\beta\,e^{(\delta\sigma^2/2)\beta^\top\bSig_0\beta}$)
gives, after differentiating in $\sigma^2$ at $0$ and using
$\beta_j(\sigma^2)=\bet_0+\delta_j\sigma^2\tilde{\bm b}_{\bSig_0}$,
exactly \eqref{app:eq:hdot}. For (a)--(b), the conditional mean of the
influence score,
\[
  \bm\psi_j(\sigma^2)
  =Y_i\bm w_i
  -\exp\!\Bigl(\beta_j(\sigma^2)^\top\bm w_i
   +\tfrac{\lambda_j\sigma^2}{2}
    \beta_j(\sigma^2)^\top\bSig_0\beta_j(\sigma^2)\Bigr)
   \bigl(\bm w_i+\lambda_j\sigma^2\bSig_0\beta_j(\sigma^2)\bigr),
\]
is real-analytic in $\sigma^2$ with
$\bm\psi_j(0)=Y_i\bm w_i-e^{\bet_0^\top\bm w_i}\bm w_i=:\bm\psi^{(0)}$
(independent of $j$), and differentiating at $\sigma^2=0$ yields
\eqref{app:eq:vdot} for
$\dot{\btV}_{jk}=\Cov(\dot{\bm\psi}_j,\bm\psi^{(0)})
+\Cov(\bm\psi^{(0)},\dot{\bm\psi}_k)
=\frac{d}{d\sigma^2}\Cov(\bm\psi_j,\bm\psi_k)\big|_{\sigma^2=0}$
(the derivative passes under $\Cov$ by dominated convergence, all
moments being uniformly bounded on $\sigma^2\le\bar\sigma^2$). Both
$\dot{\btH}$ and $\dot{\btV}$ are symmetric: for $\dot{\btV}$ this is
the general fact $\frac{d}{dt}\Cov(\bm\psi_t)
=\Cov(\dot{\bm\psi}_t,\bm\psi_t)+\Cov(\bm\psi_t,\dot{\bm\psi}_t)$.

\emph{Step 4: propagation through the sandwich.} The asymptotic
covariance of Theorem~\ref{app:thm:an} is the smooth matrix-valued map
$\bm S(\sigma^2)=\bm v_{-1}^\top\bm M\btH(\sigma^2)^{-1}
\btV(\sigma^2)\btH(\sigma^2)^{-1}\bm M^\top\bm v_{-1}$. Differentiating
$\bm H^{-1}\bm V\bm H^{-1}$ at a point where
$\bm H_0^{-1}\bm V_0\bm H_0^{-1}
=(\bm 1\bm 1^\top)\otimes\bm I_X^{-1}$ (Step~2) gives
\[
  \frac{d}{d\sigma^2}\Big|_{0}\bm H^{-1}\bm V\bm H^{-1}
  =\bm H_0^{-1}\Bigl(\dot{\btV}
    -\dot{\btH}\bm H_0^{-1}\bm V_0
    -\bm V_0\bm H_0^{-1}\dot{\btH}\Bigr)\bm H_0^{-1},
\]
whence \eqref{app:eq:cdot}. The deterministic extrapolation bias does not
enter $\bSig_{\asimex}$ (a covariance): it equals
$\sigma^2\bigl(\sum_jc_j\delta_j\bigr)\tilde{\bm b}_{\bSig_0}+O(\sigma^4)$
and affects the risk, not the variance; we account for it explicitly in
Remark~\ref{app:rem:psd} and part~(iii).

\emph{Step 5: determinant expansion.} With
$\bSig_{\asimex}=\bSig_{\mathrm{oracle}}+\sigma^2\dot{\bm C}_\beta
+O(\sigma^4)$ and both leading matrices positive definite/symmetric,
\[
  \det\bigl(\bSig_{\asimex}\bigr)
  =\det\bigl(\bSig_{\mathrm{oracle}}\bigr)
   \Bigl(1+\sigma^2\operatorname{tr}\bigl(
     \bSig_{\mathrm{oracle}}^{-1}\dot{\bm C}_\beta\bigr)
   +O(\sigma^4)\Bigr),
\]
which is part~(iii). No sign claim is made on
$\kappa_\beta$, per Remark~\ref{app:rem:psd}. This proves (ii)--(iii).

\emph{Proof of Corollary~\ref{app:cor:scalar}.} The scalar pseudo-score
condition is (Lemma~\ref{app:lem:curve}(i) and the tilt identities in its
proof)
\[
  \mu_1(\beta_0)=e^{\frac{\delta}{2}\beta^2\sigma^2}
  \bigl(\mu_1(\beta)+\delta\,\sigma^2\beta\,\mu_0(\beta)\bigr),
  \qquad \mu_r(\beta)=\E[e^{\beta X}X^r],
\]
so $\beta(\delta)=\beta_0+\delta\,\sigma^2\tilde b+O(\sigma^4)$ with
$\tilde b=-\bigl[\tfrac12\beta_0^{2}\mu_1(\beta_0)
+\beta_0\mu_0(\beta_0)\bigr]/\mu_2(\beta_0)<0$ for $\beta_0>0$; for
$X\sim N(0,1)$, $\mu_0(\beta)=e^{\beta^2/2}$,
$\mu_1(\beta)=\beta e^{\beta^2/2}$,
$\mu_2(\beta)=(1+\beta^2)e^{\beta^2/2}$, giving
$\tilde b=-\sigma^2\beta_0(1+\beta_0^2/2)/(1+\beta_0^2)$, matching
Lemma~\ref{app:lem:curve}(iii). The scalar specialization of
\eqref{app:eq:cdot} is displayed in the corollary; each ingredient
($\dot H_j$, $\dot\Gamma_{jk}$) is obtained from
\eqref{app:eq:hdot}--\eqref{app:eq:vdot} by replacing expectations with
$\mu_r(\beta_0)$-combinations (the fourth moment $\mu_4$ enters through
$\Cov(\dot\psi_j,\psi^{(0)})$, which involves
$\E[e^{2\bet_0X}X^r]$-type terms reducible to $\{\mu_r(\beta_0)\}_{r\le4}$
by completing the square). \qed
\end{proof}

